\documentclass[11pt,twoside]{article}

\usepackage[margin=1in]{geometry}
\usepackage{amsmath,amssymb,amsthm,mathtools}
\usepackage[inline,shortlabels]{enumitem}
\usepackage{xcolor}
\usepackage{../misc_files/course-definition-boxes}
\usepackage{microtype}

% Revision date: maintained by scripts/build_lecture_notes.py. Keep this command.
\newcommand{\lecturerevised}{2026-11-19 09:09 UTC-06:00}
\usepackage{eso-pic}
\AddToShipoutPictureBG{%
  \AtPageLowerLeft{%
    \put(\LenToUnit{0.5\paperwidth},\LenToUnit{0.5in}){%
      \makebox(0,0){\normalfont\footnotesize\color{black!60}Last revised: \lecturerevised}%
    }%
  }%
}

\usepackage[square,authoryear]{natbib}
\usepackage[colorlinks=true,linkcolor=blue!60!black,
  urlcolor=blue!60!black,citecolor=blue!60!black]{hyperref}
\usepackage[nameinlink,capitalize]{cleveref}
\crefname{exercise}{Exercise}{Exercises}
\Crefname{exercise}{Exercise}{Exercises}

% When including shared exercise chunks in the repository, uncomment:
\usepackage{../exercises/course-exercises}
\usepackage{needspace}

% Short refinements that may be skipped on a first reading.
\usepackage{enotez}
\DeclareInstance{enotez-list}{lecture}{list}{format=\normalsize,number={#1.}}
\setenotez{list-name=Endnotes,list-style=lecture,backref=true}
\crefname{endnote}{Endnote}{Endnotes}
\Crefname{endnote}{Endnote}{Endnotes}


\definecolor{courseblue}{HTML}{17365D}
\newcommand{\E}{\mathbb E}
\newcommand{\I}{\mathbb I}
\renewcommand{\P}{\mathbb P}
\newcommand{\R}{\mathbb R}
\newcommand{\KL}{D}
\newcommand{\TV}{\mathrm{TV}}
\makeatletter
\@for\@letter:=A,B,C,D,E,F,G,H,I,J,K,L,M,N,O,P,Q,R,S,T,U,V,W,X,Y,Z\do{%
  \expandafter\edef\csname c\@letter\endcsname{\noexpand\mathcal{\@letter}}}
\makeatother
\newcommand{\lecturenumber}{20}

\newtheorem{theorem}{Theorem}[section]
\newtheorem{proposition}[theorem]{Proposition}
\newtheorem{corollary}[theorem]{Corollary}
\theoremstyle{definition}
\newtheorem{definition}[theorem]{Definition}
\newtheorem{example}[theorem]{Example}
\theoremstyle{remark}
\newtheorem{remark}[theorem]{Remark}

\renewcommand{\thepage}{\lecturenumber--\arabic{page}}
\renewcommand{\thesection}{\lecturenumber.\arabic{section}}
\renewcommand{\theequation}{\lecturenumber.\arabic{equation}}

\pagestyle{myheadings}
\markboth{Lecture \lecturenumber: Learning from Rewards with Policy Gradients}
{Lecture \lecturenumber: Learning from Rewards with Policy Gradients}

% Course-wide notation convention:
% See learning-notation-conventions.tex: \ell is one-observation loss;
% L is population risk and L_n is empirical average loss. Keep the standard
% term empirical risk minimization (ERM). Parameterization does not change
% this convention. State local exceptions and all normalization factors.
% Every vector is a column vector. If a matrix stores one vector per position,
% row i contains the transpose of that vector. Write every row-shaped vector
% explicitly as the transpose of a column vector.
% Use n for the number of independent training examples. Use sentence case in
% optional theorem titles. Use example environments for worked examples that
% are referenced later. Reserve \boxed{} for literal boxes in diagrams.

\hypersetup{pdftitle={Lecture 20: Learning from Rewards with Policy Gradients},
  pdfauthor={Csaba Szepesvari}}
\begin{document}
\raggedbottom
\begin{center}
\fcolorbox{courseblue}{white}{%
\begin{minipage}{0.92\textwidth}
\vspace{1mm}
\textbf{CMPUT 654: Mathematical Foundations of Modern AI Systems}
\hfill Fall 2026
\vspace{3mm}
\begin{center}
{\Large\bfseries Lecture \lecturenumber: Learning from Rewards\\[1mm]
with Policy Gradients}
\vspace{1mm}

November 17, 2026
\end{center}
\vspace{2mm}
\textit{Lecturer: Csaba Szepesv\'ari}
\hfill
\textit{Notes prepared by: Csaba Szepesv\'ari}
\vspace{1mm}
\end{minipage}}
\end{center}

Reinforcement learning is a major part of training frontier models and
improving their problem-solving abilities. Policy-gradient methods are
central to this training. We study how these methods work and what
determines whether---and how quickly---they learn.

\noindent\textbf{Reading guide.}
The prerequisites are log probabilities and autoregressive factorization,
logistic classification from \href{lecture10.pdf}{Lecture 10},
gradient methods from \href{lecture11.pdf}{Lecture 11}, and the policy and
feedback terminology of \href{lecture18.pdf}{Lecture 18}.
\href{lecture19.pdf}{Lecture 19} supplies the learned-reward and
KL-regularization comparison. We first isolate sampling effects with separate
policy parameters, then use population classification problems to examine
parameter sharing. The linear examples distinguish properties of the update
from properties of the policy class; they do not give convergence guarantees
for a shared neural policy. We also examine when fitting successful responses
improves reward, and how missing contexts can undo correct behavior.
The optional section and its exercises study
changes to context weighting. The appendices give supplementary proofs and
extensions, including the online and natural-gradient counterexamples.

\section{From rewards to parameter updates}
\label{sec:feedback}

Consider training a language model to solve programming problems. Starting
from a pretrained model, we generate candidate programs, run tests, and use
the resulting rewards to update the model. The aim is to increase the
probability that future programs pass the tests. A test supplies an
evaluation, while a parameter update requires a direction in parameter
space. Our first task is to derive such a direction from the generated
program and its reward.

Lecture 17 used verification to select outputs from a fixed generator.
Here we use evaluations to train the generator. Lecture 18 established
advantages of collecting data interactively; the question now is how
policy-gradient updates make use of those data.

Language-model training can also involve genuine multi-turn interaction.
For example, a model may issue a search query or a program call, observe the
result, and then choose its next action; a user may also respond between model
turns. An earlier model action then changes the information available later.
The policy therefore determines a distribution over complete interaction
trajectories, and a terminal task reward must be assigned to several model
actions. This setting is a sequential reinforcement-learning problem.
Autoregressive generation of one response also contains many token decisions.
When no external observation arrives before the response ends, those decisions
can be treated as the single sampled response $Y$ below. Multi-turn interaction
requires the sequential formulation because later model decisions condition on
new external observations. \Cref{app:return} gives the trajectory estimator,
returns to go, and history-dependent value baselines.

Fix a prompt $x$ and a finite set $\cY$ of complete responses. The prompt
was denoted by $u$ in Lectures 18--19; here $x$ also serves as the input
in our classification examples, and $u$ will be a scalar policy parameter.
Let $\theta\in\R^d$ parameterize a differentiable, strictly positive policy
$\pi_\theta(\cdot\mid x)$. Generate $Y\sim\pi_\theta(\cdot\mid x)$,
then observe a bounded reward $R$. The conditional reward law given $Y=y$
is fixed independently of $\theta$, with mean $r(x,y)$. The objective is
\begin{equation}
 J_x(\theta)=\E[R]=\sum_{y\in\cY}\pi_\theta(y\mid x)r(x,y).
 \label{eq:reward-objective}
\end{equation}
The symbol $J_x$ denotes expected reward, which we maximize. This is
Lecture 18's policy objective at a fixed prompt, expressed as a function
of the policy parameter $\theta$. Gradients are with respect to $\theta$
unless another variable is indicated.
A deterministic verifier is the special case $R=r(x,Y)\in\{0,1\}$.
Learning with such computable feedback is called \emph{reinforcement learning
with verifiable feedback} (RLVF), also called \emph{reinforcement learning with
verifiable rewards} (RLVR). A learned reward model supplies another
possible evaluator. In either case, $J_x$ measures the criterion that the
evaluator supplies.

\subsection{A derivative that can be sampled}

Enumerating \cref{eq:reward-objective} may be impossible, and a test need not
be differentiable. We therefore seek an expression for $\nabla J_x$ using
one generated response, its reward, and derivatives of the policy.
Differentiating the finite sum gives
$\nabla J_x=\sum_y r(x,y)\nabla\pi_\theta(y\mid x)$.
The sampler supplies weights $\pi_\theta(y\mid x)$, so we insert those
weights by writing
\[
 \nabla\pi_\theta(y\mid x)
 =\pi_\theta(y\mid x)\nabla\log\pi_\theta(y\mid x).
\]
The \emph{score} of a prompt--response pair is the vector
$s_\theta(x,y)=\nabla\log\pi_\theta(y\mid x)$.
The logarithm appears because its derivative converts the available
sampling distribution into the derivative we want.

\begin{proposition}[Score identity and baseline cancellation]
\label{prop:score}
Under the finite-response setup above,
\begin{equation}
 \nabla J_x(\theta)=\E[R s_\theta(x,Y)],
 \qquad \E[s_\theta(x,Y)]=0.
 \label{eq:score-identities}
\end{equation}
For every fixed scalar $b$,
\begin{equation}
 \E[(R-b)s_\theta(x,Y)]=\nabla J_x(\theta).
 \label{eq:baseline-identity}
\end{equation}
\end{proposition}
\begin{proof}
Conditioning on the sampled response replaces its reward by its mean:
\[
 \E[R s_\theta(x,Y)]
 =\sum_y\pi_\theta(y\mid x)r(x,y)s_\theta(x,y)
 =\sum_y r(x,y)\nabla\pi_\theta(y\mid x)=\nabla J_x(\theta).
\]
The score has mean zero because policy probabilities sum to one:
\[
 \E[s_\theta(x,Y)]
 =\sum_y\nabla\pi_\theta(y\mid x)
 =\nabla\!\left(\sum_y\pi_\theta(y\mid x)\right)=0.
\]
Subtracting $b\E[s_\theta(x,Y)]$ proves
\cref{eq:baseline-identity}.
\end{proof}

A \emph{baseline} is the reference value $b$ subtracted from the reward in
the score estimator.
The baseline can depend on the prompt, current parameters, and
previous observations, provided it is fixed before the current response is
sampled. Conditional on that history, the same proof applies.
A baseline using the current response requires its own calculation;
\cref{sec:groups} studies that case. When forming the estimator, we hold
$b$ fixed, even if its numerical value was computed from $\theta$.
Differentiating through $b$ would define a different update.

The policy's \emph{value} at prompt $x$ is its expected reward
$J_x(\theta)$. The \emph{advantage} of response $y$ is
$r(x,y)-J_x(\theta)$, its expected reward above the policy's average at
that prompt. With the exact value as baseline,
\[
 \E[R-J_x(\theta)\mid Y=y]=r(x,y)-J_x(\theta).
\]
When $b$ estimates the value, $R-b$ is used as an \emph{advantage estimate}.
One can fit a value predictor $v_w(x)$ to rewards sampled from the current
policy and use its predictions as baselines for later responses.
Lecture 11's temporal-difference example uses $v_w(s)$ for a state-value
predictor; \cref{app:return} gives the sequential version of the advantage.
One motivation for subtracting a baseline is to reduce the variance of
the gradient estimator. The value baseline need not minimize that variance;
\cref{ex:l20-variance-baseline} asks you to find the minimizing baseline.

\begin{remark}[Other unbiased estimators]
\label{rem:score-characterization}
One may ask whether there are other unbiased estimators of the gradient
from one evaluated response. Essentially, the answer is no: if the same
estimator must work for every unknown reward function, it has the form
$R s_\theta(x,Y)+h(Y)$, where $h$ is independent of the reward
function and $\E[h(Y)]=0$. Scalar baselines are one choice of this
mean-zero correction.\endnote{Fix $\theta,x$ and write
$p_y=\pi_\theta(y\mid x)$, $g_y=\nabla\pi_\theta(y\mid x)$,
and $h_y=h(y)$. A scalar baseline chooses $h_y=-b g_y/p_y$.
For any known vectors $c_y$, the choice
$h_y=c_y-\sum_jp_jc_j$ has mean zero. A known reward predictor
$\widehat r_y$ gives the unbiased estimator
$(R-\widehat r_Y)s_\theta(x,Y)+\sum_y\widehat r_y g_y$;
computing the correction may require enumerating responses. The
characterization permits arbitrary reward functions. For two actions with
exactly one correct action and deterministic correctness reward,
$(2R-1)\nabla\pi_\theta(Y\mid x)$ always equals the gradient of the
correct-action probability: a failure identifies the other action, whose
probability gradient is $-\nabla\pi_\theta(Y\mid x)$. This zero-variance
estimator uses the restricted reward family. For the universal-reward
requirement, policy probabilities are needed: access to $\nabla\pi$ alone,
without $\pi$, does not generally
determine the required ratio.}
\Cref{ex:l20-unbiased-estimators} asks you to characterize all such estimators.
\end{remark}

At iteration $t=0,1,\ldots$, draw $Y_t$ from
$\pi_{\theta_t}(\cdot\mid x)$ using fresh randomness and observe $R_t$.
Throughout the lecture, parameter iterates and quantities computed from
them retain lowercase notation, even when random; sampled observations
use capitals. Using a scalar baseline $b_t$ chosen before drawing $Y_t$
and step size $\eta_t>0$ gives the
\emph{stochastic policy-gradient} update
\begin{equation}
 \theta_{t+1}=\theta_t+
 \eta_t(R_t-b_t)s_{\theta_t}(Y_t\mid x).
 \label{eq:sampled-update}
\end{equation}
This score-function update is known as \emph{REINFORCE}.
One update requires generating a response, evaluating it once, and computing
a gradient of its log probability. A batch averages several such directions.
Generation, evaluation, and the backward computation all contribute to cost.
For prompts sampled from a fixed distribution on a finite prompt set,
independent of $\theta$,
averaging \cref{eq:score-identities} over the prompt gives the gradient of
average reward across prompts.

\subsection{Token scores and complete responses}

The response-level estimator can be implemented with an autoregressive
model. Write $p_\theta(\cdot\mid x,y_{<h})$ for the next-token model,
as in Lectures 1 and 19. For a fixed generation length $H$, write a
response as $y=y_{1:H}$ and its prefix as $y_{<h}=y_{1:h-1}$, with
$y_{<1}$ the empty sequence. Then
\begin{equation}
 \pi_\theta(y_{1:H}\mid x)=\prod_{h=1}^H
 p_\theta(y_h\mid x,y_{<h}),\qquad
 s_\theta(x,y_{1:H})=\sum_{h=1}^H
 \nabla\log p_\theta(y_h\mid x,y_{<h}).
 \label{eq:token-score}
\end{equation}
\endnote{For generation that stops at a sampled end-of-sequence token
$\mathrm{EOS}$, assume a maximum length $H$ and pad after EOS with an
absorbing symbol of conditional probability one, as in Lecture 18.
The same fixed-$H$ formula applies to the padded responses. Sampled EOS
contributes its score; padding contributes zero. Equivalently, without
padding, the product and sum run to the actual length $|y|$, including
EOS when sampled, as in Lecture 19. Truncation at $H$ adds no sampled
decision and hence no score term. The EOS factors already encode the
probability of termination.}
A terminal reward multiplies the sum
of token scores; computing the sum does not require differentiating the
verifier. The update changes token probabilities through the shared model
parameters. It supplies no separate observation identifying which token
caused success or failure.

\paragraph{Intermediate rewards.}
Suppose feedback arrives after individual tokens, with $R_{h+1}$ observed
after $Y_h$, and the objective is the expected total reward
$J(\theta)=\E[R]$, where $R=\sum_{h=1}^H R_{h+1}$.
Evaluations of intermediate reasoning steps, often called \emph{process
feedback}, can supply such rewards. Assume that the feedback law given the
generated prefix is independent of $\theta$. Writing
$s_h=\nabla\log p_\theta(Y_h\mid x,Y_{<h})$, we can replace the full return
by the \emph{return to go} in each score term:
\[
 \nabla J(\theta)
 =\E\!\left[\sum_{h=1}^H s_h R\right]
 =\E\!\left[\sum_{h=1}^H s_h\sum_{k=h}^H R_{k+1}\right].
\]
Past rewards cannot be affected by the current token, and their score terms
have expectation zero. Removing them is intended to reduce variance;
the variance of the full sum need not decrease in every problem.\endnote{For
example, take two independent Bernoulli actions with separate logits and
constant rewards $R_2=1$ and $R_3=-1$. The full-return estimator is identically
zero, whereas the return-to-go estimator is $-s_2$ and has positive variance.
The baseline calculation in \cref{ex:l20-variance-baseline} provides another
way to see why subtracting a mean-zero score term need not reduce variance.}
\Cref{app:return} proves the identity and explains its connection to
Rao--Blackwellization.

\section{Softmax bandits}
\label{sec:bandits}

We now study whether sampled policy-gradient updates learn and how long
learning takes. Fix one prompt and give each possible response its own
logit. The resulting model is a \emph{multi-armed bandit}: the learner
chooses an action and observes a reward for that action. Separate logits
remove parameter sharing while retaining the dependence of future samples
on the current policy. We will establish eventual convergence, prove a
learning-time bound for deterministic binary feedback, and then see how
stochastic rewards can make learning much slower. We then identify a small
fixed-step regime that recovers the $1/t$ expected rate.

\subsection{Setup and update}
\label{sec:bandit-setup}

Let $\cA=[K]$ be the action set, with $K\ge2$. Each action $a$ has an
unknown, fixed \emph{reward distribution} $\mu_a$ on $[-B,B]$, where
$0<B<\infty$, and mean reward
\[
 r_a=\int z\,\mu_a(dz).
\]
For a logit vector $\theta\in\R^K$, use the softmax policy and objective
\[
 \pi_\theta(a)=\frac{e^{\theta_a}}{\sum_{j=1}^K e^{\theta_j}},
 \qquad J(\theta)=\sum_{a=1}^K\pi_\theta(a)r_a.
\]
The fixed prompt is suppressed. Write $a^\star\in\arg\max_a r_a$ and
\[
 J^\star=\sup_{\theta\in\R^K}J(\theta)=\max_a r_a=r_{a^\star}.
\]
The equality follows by concentrating softmax probability on a best
action. The supremum need not be attained at a finite parameter.

Start from a finite deterministic $\theta_0$ and write
$\pi_t=\pi_{\theta_t}$. Let $\cF_t$ be the history generated by
$(A_0,R_0),\ldots,(A_{t-1},R_{t-1})$, before round $t$. At that round,
draw $A_t$ from $\pi_t$ and obtain a fresh reward with distribution
$\mu_{A_t}$. Precisely, for every $a\in\cA$ and measurable set $C\subseteq\R$,
\begin{equation}
 \P(A_t=a,\ R_t\in C\mid\cF_t)=\pi_t(a)\mu_a(C).
 \label{eq:bandit-sampling}
\end{equation}
Thus the action probabilities depend on the history, while the reward law
given the action remains fixed. In particular,
$\E[R_t\mid\cF_t]=J(\theta_t)$.

\begin{remark}[Sampled softmax update]
\label{rem:sampled-softmax}
For the separate-logit softmax policy, viewed as a probability vector,
\[
 \nabla\log\pi_\theta(a)=e_a-\pi_\theta,
\]
where $e_a$ is the $a$th standard basis vector. With zero baseline and a
fixed step size $\eta>0$, \cref{eq:sampled-update} therefore becomes
\begin{equation}
 \theta_{t+1}=\theta_t+\eta R_t(e_{A_t}-\pi_t).
 \label{eq:vector-pg}
\end{equation}
Verify the score identity and derive the update from \cref{eq:sampled-update}.
\end{remark}

We measure both the quality of the current policy and the reward lost
during training. Define the \emph{suboptimality gap} and
\emph{expected cumulative regret}
by
\[
 \delta_t=J^\star-J(\theta_t),\qquad
 \cR_T=\sum_{t=0}^{T-1}\E[\delta_t]
       =TJ^\star-\E\!\left[\sum_{t=0}^{T-1}R_t\right].
\]
The last equality follows from $\E[R_t\mid\cF_t]=J(\theta_t)$.
The benchmark $TJ^\star$ is the expected reward from always choosing a
best action. Thus $\cR_T/T$ is the expected reward lost per training round.

\paragraph{An exact-gradient comparison.}
Before studying the sampled update, consider whether a large suboptimality gap
must be visible in the exact gradient.

\begin{proposition}[A probability-dependent gradient bound]
\label{prop:gradient-gap}
For the separate-logit softmax policy,
\begin{equation}
 \partial_{\theta_a}J(\theta)=\pi_\theta(a)(r_a-J(\theta)),
 \qquad
 \pi_\theta(a^\star)\bigl(J^\star-J(\theta)\bigr)
 \le\|\nabla J(\theta)\|_2.
 \label{eq:gradient-gap}
\end{equation}
\end{proposition}
\begin{proof}
The softmax derivative is
$\partial_{\theta_a}\pi_\theta(j)=\pi_\theta(j)
(\I\{j=a\}-\pi_\theta(a))$.
Substitution into $\partial_{\theta_a}J$ gives the first identity.
The coordinate for $a^\star$ is nonnegative and at most the Euclidean norm
of the gradient, proving the inequality.
\end{proof}

The bound connects to gradient domination in Lecture 11:
\[
 J^\star-J(\theta)\le
 \frac{\|\nabla J(\theta)\|_2}{\pi_\theta(a^\star)}.
\]
A small gradient certifies a small gap only when the optimal action has
enough probability. The same probability also determines how often a
sampled update observes that action. The convergence results below must
account for this dependence on the current policy.

\subsection{Almost-sure convergence}
\label{sec:time}

For the separate-logit update, every positive constant step size eventually
learns a best action under the following assumptions. The statement concerns
the entire sampled training trajectory.

\begin{theorem}[Constant-step almost-sure convergence]
\label{thm:as}
In the setup of \cref{sec:bandit-setup}, suppose the mean rewards $r_a$
are distinct. For every fixed $\eta>0$, the update \cref{eq:vector-pg}
satisfies
\[
 \pi_t(a^\star)\longrightarrow1,
 \qquad \delta_t\longrightarrow0
 \qquad\text{almost surely}.
\]
\end{theorem}

The proof is outside of the scope of this lecture. The theorem permits arbitrary positive
constant steps. Bounded rewards also give a consequence for expected regret.

\begin{corollary}[Sublinear expected regret]
\label{cor:sublinear-regret}
Under the assumptions of \cref{thm:as},
\[
 \frac{\cR_T}{T}\longrightarrow0\qquad(T\to\infty).
\]
\end{corollary}

\Cref{ex:l20-sublinear-regret} asks you to prove the corollary.
\emph{Hint:} First show that $\E[\delta_t]\to0$.
Neither statement gives a learning-time bound. We next obtain one for a
deterministic binary verifier.

\subsection{A learning-time bound for a binary verifier}
\label{sec:binary}

How many generated and evaluated responses are sufficient to make failures
uncommon? We answer this question for a policy choosing between two options
at a fixed input $x$. For example, $x$ may encode a question and a proposed
solution, and the policy must decide whether to accept ($a=1$) or reject
($a=0$) the solution. Assume the reward is deterministically one for a
correct decision and zero otherwise. Without loss of generality, suppose
$r_a=a$; the other case follows by exchanging the action labels.

The reward laws on these actions satisfy
$\mu_0(\{0\})=\mu_1(\{1\})=1$, so $R_t=A_t$.
For the two-logit policy, put
\[
 u_t=\theta_{t,1}-\theta_{t,0},\qquad
 p_t=\pi_t(1)=\sigma(u_t)=\frac1{1+e^{-u_t}},\qquad q_t=1-p_t.
\]
Subtracting the two coordinate updates in \cref{eq:vector-pg} gives
\begin{equation}
 u_{t+1}=u_t+2\eta A_t(1-p_t).
 \label{eq:binary-update}
\end{equation}
Both logits move, which accounts for the factor $2$.
\Cref{ex:l20-binary-update} asks you to derive the reduction.
Conditional on $\cF_t$, $A_t$ is Bernoulli with parameter $p_t$.
Here $J^\star=1$ and $\delta_t=q_t$. Put $m_t=\E[q_t]$; the expectation
averages the trained policy's failure probability over its training samples.

\begin{theorem}[Learning from a binary verifier]
\label{thm:binary}
For \cref{eq:binary-update}, with finite deterministic $u_0$ and any fixed
$\eta>0$, let $c_0=2\eta e^{-2\eta}p_0^2$. After $t\ge0$ updates,
\begin{equation}
 \E[\delta_t]=\E[1-p_t]
 \le\frac{1}{\delta_0^{-1}+c_0t}.
 \label{eq:binary-rate}
\end{equation}
Moreover, $p_t\to1$ almost surely.
\end{theorem}

\begin{corollary}[Logarithmic regret for a binary verifier]
\label{cor:binary-regret}
Under the assumptions of \cref{thm:binary}, for every integer $T\ge1$,
\[
 \cR_T\le q_0+\frac1{c_0}\log\!\bigl(1+c_0q_0(T-1)\bigr).
\]
In particular, $\cR_T=O(\log T)$ for fixed $\eta$ and $p_0$ as $T\to\infty$.
\end{corollary}
\begin{proof}
Here $\delta_t=q_t$. Sum \cref{eq:binary-rate} and compare the decreasing
summands after $t=0$ with an integral:
\[
 \cR_T\le\sum_{t=0}^{T-1}\frac1{q_0^{-1}+c_0t}
 \le q_0+\int_0^{T-1}\frac{ds}{q_0^{-1}+c_0s}
 =q_0+\frac1{c_0}\log\!\bigl(1+c_0q_0(T-1)\bigr).
\]
\end{proof}

To prove \cref{thm:binary}, we first quantify what one successful response
does to the policy, then accumulate the expected progress over training.

\subsubsection{Finding the progress estimate}

A failure gives no update. A success increases $u_t$ by $2\eta q_t$.
Consequently, $p_t$ never decreases and $p_0\le p_t$ on every trajectory.
Before the first success the policy stays fixed, so the number of evaluations
up to and including that success has mean $1/p_0$. A very small initial
success probability already imposes a long wait.

To predict the progress after a success, write $p=p_t$ and $q=q_t$.
The derivative $dp/du=pq$ suggests that a small successful logit change
$2\eta q$ increases $p$ by approximately
\[
 (pq)(2\eta q)=2\eta p q^2.
\]
A success occurs with probability $p$, so expected progress should be
approximately $2\eta p^2q^2$. The two factors of $q$ have distinct sources:
the logit update shrinks, and the logistic policy becomes less sensitive to
logit changes as it approaches probability one. We therefore look for a
finite-step lower bound with this dependence. The approximation motivates
the bound; it will not be used as a proof step.

Let $q^+$ denote the failure probability after a successful update from
$(p,q)$. Success multiplies the success-to-failure odds by $e^{2\eta q}$.
Converting the new odds back to a probability gives
\begin{equation}
 q^+=\frac{q}{q+pe^{2\eta q}},\qquad
 q-q^+=\frac{pq(e^{2\eta q}-1)}{q+pe^{2\eta q}}.
 \label{eq:finite-progress}
\end{equation}
We want a lower bound on improvement. Lower-bound the numerator and
upper-bound the denominator. Convexity of the exponential gives
$2\eta q\le e^{2\eta q}-1$, while $p+q=1$ and $q\le1$ give
\[
 q+pe^{2\eta q}\le e^{2\eta q}\le e^{2\eta}.
\]
The two elementary bounds yield
\begin{equation}
 2\eta e^{-2\eta}pq^2\le q-q^+.
 \label{eq:success-progress}
\end{equation}
The finite-step estimate retains the predicted dependence, with the
additional factor $e^{-2\eta}$. Multiplying by the probability $p_t$ of the
successful observation, and using $p_0\le p_t$, gives
\begin{equation}
 c_0 q_t^2\le\E[q_t-q_{t+1}\mid\cF_t],
 \qquad c_0=2\eta e^{-2\eta}p_0^2>0.
 \label{eq:binary-drift}
\end{equation}
The constant $c_0$ is deterministic. We can therefore take expectations and
use $m_t^2\le\E[q_t^2]$ to close the recurrence:
\begin{equation}
 m_{t+1}\le m_t-c_0m_t^2.
 \label{eq:quadratic-drift}
\end{equation}

\subsubsection{Accumulating progress}

The policy-gradient calculation has produced a scalar recurrence. Extracting
a rate from that recurrence is a standard, reusable argument; the same
quadratic recurrence appeared in Lecture 11's analysis of separable
cross-entropy. We now explain the rate before proving the general bound.

Suppose an error is of size $a$ and decreases by approximately $ca^2$ per
step. Halving the error requires a change of size $a/2$, so the number of
steps should be of order $a/(ca^2)=1/(ca)$. At the next, smaller error
scale, the required change halves but the progress per step quarters.
Successive halvings therefore take about twice as long. Adding those times
gives a total dominated by the last halving: reaching error $\varepsilon$
should cost order $1/(c\varepsilon)$ steps. This predicts the rate when
progress has the stated size; a lower bound on progress will give an upper
bound on the required time.

How can we make that calculation precise without splitting time into
halvings? Divide each small decrease in error by the progress available at
that error. For quadratic progress, accumulating these ratios gives
\[
 \int_\varepsilon^{a_0}\frac{du}{cu^2}
 =\frac1c\left(\frac1\varepsilon-\frac1{a_0}\right).
\]
The reciprocal appears because it is the integral of inverse quadratic
progress. The discrete proof compares this integral with one rectangle
per step. The same construction works for other powers.

\begin{proposition}[Standard drift-to-rate bound]
\label{prop:drift-rate}
Let $(a_t)_{t\ge0}$ be a nonnegative sequence with $a_0>0$, and suppose
\[
 a_{t+1}\le a_t-ca_t^\alpha
 \qquad(t\ge0)
\]
for constants $c>0$ and $\alpha>1$. Then
\[
 a_t\le\bigl(a_0^{1-\alpha}+c(\alpha-1)t\bigr)^{-1/(\alpha-1)}.
\]
In particular, quadratic decrease ($\alpha=2$) gives
$a_t\le(a_0^{-1}+ct)^{-1}$.
\end{proposition}
\begin{proof}
The recurrence makes $a_t$ nonincreasing. If an iterate reaches zero,
all subsequent iterates are zero, so consider steps with $a_{t+1}>0$.
One step decreases the error by at least $ca_t^\alpha$.
On the interval $[a_{t+1},a_t]$, the function $u^{-\alpha}$ is at
least $a_t^{-\alpha}$. Therefore
\[
 c\le\frac{a_t-a_{t+1}}{a_t^\alpha}
 \le\int_{a_{t+1}}^{a_t}u^{-\alpha}\,du.
\]
The intervals from successive steps are adjacent. Summing their integrals
from step zero through step $t-1$ gives
\[
 ct\le\int_{a_t}^{a_0}u^{-\alpha}\,du
 =\frac{a_t^{1-\alpha}-a_0^{1-\alpha}}{\alpha-1}.
\]
Rearranging proves the bound. For quadratic progress, each integral is
exactly $1/a_{t+1}-1/a_t$: the familiar reciprocal proof is the same
calculation written without integrals.
\end{proof}

The reusable part needs only a deterministic recurrence. For our sampled
algorithm, the separate tasks were to prove conditional drift, take
expectations, and use convexity to obtain that recurrence for $m_t$.


\begin{proof}[Completion of the proof of \cref{thm:binary}]
Apply \cref{prop:drift-rate} to \cref{eq:quadratic-drift} with
$a_t=m_t$, $a_0=q_0$, $c=c_0$, and $\alpha=2$ to obtain
\cref{eq:binary-rate}.
For almost-sure convergence, the decreasing sequence $q_t$ has a
nonnegative limit $q_\infty$. Bounded convergence and the rate give
$\E[q_\infty]=\lim_t\E[q_t]=0$, so $q_\infty=0$ almost surely.
\end{proof}

\Cref{ex:l20-binary-baseline} asks you to extend the theorem to any fixed
baseline $b\in[0,1]$, with the same rate.

The bound gives a sufficient evaluation budget: for $0<\varepsilon<q_0$,
\[
 t\ge\frac{\varepsilon^{-1}-q_0^{-1}}{2\eta e^{-2\eta}p_0^2}
 \quad\Longrightarrow\quad \E[q_t]\le\varepsilon.
\]
The dependence on initialization need not be sharp; the proof replaced every
$p_t$ by $p_0$. \Cref{ex:l20-first-success} separates a necessary waiting
cost from this sufficient guarantee. Likewise, the factor $2\eta e^{-2\eta}$
comes from a coarse finite-step estimate and does not identify an optimal
step size.

The proof uses deterministic zero/one rewards: an incorrect response gives
no update, and a correct response increases the logit difference. Thus the
success probability never decreases. We next change the reward laws while
keeping the same softmax update and see how this monotonicity can fail.

\subsection{Nearly linear regret with stochastic rewards}
\label{sec:slow}

Consider two actions with $\theta_0=0$ and reward laws
\[
 \mu_1(\{+1\})=\frac{1+\Delta}{2},\qquad
 \mu_1(\{-1\})=\frac{1-\Delta}{2},\qquad
 \mu_2(\{0\})=1,\qquad 0<\Delta<1.
\]
The first law is a biased Rademacher distribution. Its mean is
$r_1=\Delta$, while $r_2=0$, so action 1 is optimal.

A sequence of negative rewards from action 1 reduces its sampling
probability. Once that probability is at most $1/(2T)$, repeatedly choosing
action 2 leaves the logits unchanged because its reward is zero. For the
next $T$ evaluations, the conditional probability of never choosing action 1
is then at least
\[
 \left(1-\frac1{2T}\right)^T\ge\frac12.
\]
Thus reaching a very small optimal-action probability can cause a long period
with little corrective information. A regret lower bound must also quantify
the probability of reaching such a state. Balancing that probability against
the cost of the subsequent delay yields the following result.

\begin{theorem}[Polynomial expected regret]
\label{thm:slow}
For the two-action instance above and update \cref{eq:vector-pg}, define
$\lambda_\Delta=\tfrac12\log((1+\Delta)/(1-\Delta))$.
If $\lambda_\Delta<\eta$, then, for every $\varepsilon>0$,
\begin{equation}
 \cR_T=\widetilde\Omega\!\left(
 T^{1-(1+\varepsilon)\lambda_\Delta/\eta}\right)
 \qquad(T\to\infty).
 \label{eq:slow-regret}
\end{equation}
Here the notation suppresses logarithmic factors in $T$ and constants
depending on the fixed instance, step size, and $\varepsilon$.
\end{theorem}

The theorem is stated without its technical lower-bound proof. For any fixed
$\eta>0$ and $\alpha\in(0,1)$, choose a fixed $\Delta>0$ sufficiently small
that $\lambda_\Delta<\alpha\eta/2$. Taking $\varepsilon=1$ in the theorem gives
\[
 \cR_T=\widetilde\Omega(T^{1-\alpha}).
\]
Thus the regret exponent can be arbitrarily close to one across fixed
instances, even for a fixed step size.

The reward laws satisfy the assumptions of \cref{thm:as}, so this same
instance has $\pi_t(1)\to1$ almost surely and $\cR_T/T\to0$.
A polynomial lower bound with exponent below one is compatible with
sublinear regret. The convergence theorem places no quantitative bound on
how long unsuccessful training trajectories can last.

The logarithmic bound in \cref{cor:binary-regret} used deterministic
zero/one rewards. Here a negative reward from the optimal action reduces
its probability, while a draw of the zero-reward action leaves the policy
unchanged. The resulting delays invalidate the monotonicity argument and
can make expected regret polynomial.

\subsection{Small fixed steps recover the \texorpdfstring{$1/t$}{1/t} rate}
\label{sec:fixed-step-rate}

The preceding lower bound leaves open whether stochastic rewards themselves
prevent a $1/t$ expected rate. A sufficiently small, horizon-independent
fixed step recovers that rate. The theorem permits signed rewards and an
arbitrary finite deterministic initialization, but requires every pair of
arms to have distinct mean rewards.

Define the smallest and largest gaps from the unique best arm and the
smallest pairwise mean separation by
\[
 \Delta_{\min}=\min_{a\ne a^\star}(J^\star-r_a),\qquad
 \Delta_{\max}=\max_{a\ne a^\star}(J^\star-r_a),\qquad
 \Delta_{\mathrm{pair}}=\min_{a\ne b}|r_a-r_b|.
\]
The stepsize below depends on $\Delta_{\mathrm{pair}}$, which also compares
pairs of suboptimal arms, rather than only on $\Delta_{\min}$.

\begin{theorem}[Expected rate for a small fixed step]
\label{thm:fixed-step-rate}
In the setup of \cref{sec:bandit-setup}, suppose the mean rewards are
distinct. Define
\begin{equation}
 M=\prod_{a\ne a^\star}
 \left(1+\frac{3\Delta_{\max}}{\Delta_{\min}}
 \exp(\theta_{0,a}-\theta_{0,a^\star})\right)^2.
 \label{eq:fixed-step-constant}
\end{equation}
If
\begin{equation}
 0<\eta\le\frac{\Delta_{\mathrm{pair}}^2}{40\sqrt K B^3},
 \label{eq:fixed-step-condition}
\end{equation}
then, for every $t\ge0$,
\begin{equation}
 \E[\delta_t]
 \le\frac{\delta_0}
 {1+\dfrac{\eta\delta_0}{2M}t}.
 \label{eq:fixed-step-rate}
\end{equation}
In particular, for $t\ge1$,
\[
 \E[\delta_t]\le\frac{2M}{\eta t}.
\]
Moreover, $\pi_t(a^\star)\to1$ almost surely.
\end{theorem}

\begin{corollary}[Logarithmic regret for a small fixed step]
\label{cor:fixed-step-regret}
Under the assumptions of \cref{thm:fixed-step-rate}, for every integer
$T\ge1$,
\[
 \cR_T\le \delta_0+\frac{2M}{\eta}
 \log\!\left(1+\frac{\eta\delta_0}{2M}(T-1)\right).
\]
Consequently, $\cR_T=O(\log T)$ for a fixed bandit instance, stepsize,
and initialization.
\end{corollary}

The main proof problem is that the probability of the best arm and the
current suboptimality gap depend on the same training history. The proof
first constructs a polynomial Lyapunov function that gives the uniform
inverse-moment bound
\begin{equation}
 \sup_{t\ge0}\E[\pi_t(a^\star)^{-2}]\le M.
 \label{eq:fixed-step-moment}
\end{equation}
A one-step smoothness calculation gives the conditional progress estimate
\begin{equation}
 \delta_t-\E[\delta_{t+1}\mid\cF_t]
 \ge\frac{\eta}{2}\pi_t(a^\star)^2\delta_t^2.
 \label{eq:fixed-step-progress}
\end{equation}
Put $m_t=\E[\delta_t]$. Cauchy--Schwarz combines the two estimates without
assuming independence:
\[
 m_t^2
 =\left(\E[(\pi_t(a^\star)\delta_t)\pi_t(a^\star)^{-1}]\right)^2
 \le M\E[\pi_t(a^\star)^2\delta_t^2].
\]
Taking expectations in \cref{eq:fixed-step-progress} therefore yields
\begin{equation}
 m_t-m_{t+1}\ge\frac{\eta}{2M}m_t^2.
 \label{eq:fixed-step-recurrence}
\end{equation}
The standard drift-to-rate bound in \cref{prop:drift-rate} proves
\cref{eq:fixed-step-rate}; summing that rate proves
\cref{cor:fixed-step-regret}. The polynomial Lyapunov calculation establishing
\cref{eq:fixed-step-moment} is the technical part of the proof and is omitted.

The constant $M$ can be large. Under uniform initialization,
\[
 M=\left(1+\frac{3\Delta_{\max}}{\Delta_{\min}}
 \right)^{2(K-1)},
\]
so the theorem's guarantee can be exponential in $K$. The proof does not show
that this particular exponential dependence is necessary under uniform
initialization. Severe dependence on initial coverage is nevertheless real.
The construction in \citet[Theorem~5.6]{mei2024} has a suboptimal plateau
whose length contains the factor $1/\pi_0(a^\star)$. With the other
construction parameters fixed, an exponentially small initial probability
of the best arm therefore produces an exponentially long plateau. The
displayed value of $M$ is an upper bound whose precise dependence on $K$ may
be loose; the theorem does not give a polynomial guarantee in $K$ or in
$1/\pi_0(a^\star)$.

For the two-arm instance in \cref{sec:slow}, $B=1$ and
$\Delta_{\mathrm{pair}}=\Delta$.
The sufficient condition in \cref{thm:fixed-step-rate} is then
$\eta\le\Delta^2/(40\sqrt2)$, whereas \cref{thm:slow} applies when
$\eta>\lambda_\Delta$, with $\lambda_\Delta\sim\Delta$ as
$\Delta\downarrow0$. The two theorems leave an intermediate stepsize regime
between these sufficient conditions.

\Cref{app:small-step} states a complementary finite-horizon guarantee that
allows tied mean rewards. Decreasing steps also reduce later corrections to
an early preference for a suboptimal arm. \Cref{app:diffusion} develops a
related lower bound in a continuous-time model; that result is supplementary.

\section{Learning with shared parameters}
\label{sec:sharing}

The preceding convergence guarantees use separate logits for one prompt.
A language model shares parameters across prompts: an update that improves
one response can change the probabilities of responses to other prompts.
We now examine how much this changes the convergence question. Expected
reward still measures the desired performance. We first identify a region
where its gradient guarantees learning, then show how exact policy gradient
can fail to reach that region even on a separable classification problem.
We then compare reward maximization with fitting successful responses,
separating restrictions of the policy class from missing training contexts.

\subsection{Classification as a contextual bandit}

Classification lets us compare reward maximization and supervised learning
on the same prediction problem. Let $X\sim\nu$ take values in $\R^d$,
with $\E[\|X\|_2^2]<\infty$, and let $y:\R^d\to\{-1,1\}$ be the
measurable correct-label function. Given $X=x$, draw a label $A\in\{-1,1\}$ from
the linear logistic policy
\[
 \pi_\theta(a\mid x)
 =\frac{\exp(a\theta^\top x/2)}
 {\exp(-\theta^\top x/2)+\exp(\theta^\top x/2)}
 =\sigma(a\theta^\top x).
\]
This is the ordinary two-action softmax policy. The two logits are
$-\theta^\top x/2$ and $\theta^\top x/2$; subtracting a common logit gives
the equivalent pair $(0,\theta^\top x)$ used in the two-action reduction of
\cref{sec:binary}. The same policy is Lecture 10's binary predictor $q_w$
with $w=\theta$, now sampled to choose an action. The log-odds
$\theta^\top x$ is the quantity denoted there by $s_w(x)$. The policy score
$s_\theta(x,a)=\nabla_\theta\log\pi_\theta(a\mid x)$ is a vector.
Define $r(x,a)=\I\{a=y(x)\}$ and observe $R=r(X,A)$.
This is a \emph{contextual bandit}:
the learner observes the context, chooses one action, and receives its
reward. In this model the logits are linear in the features; reward
linearity is a separate condition.

Following Lectures 10--11, write $z(x)=y(x)x$ for the
\emph{signed feature vector}. Its inner product
$m_\theta(x)=z(x)^\top\theta$ is the \emph{signed margin}, also the
log odds of choosing the correct label. The probability of choosing that
label is $p_\theta(x)=\pi_\theta(y(x)\mid x)=\sigma(m_\theta(x))$.

\paragraph{Reward and log loss.}
The log loss at $x$ is
$\ell(\theta,x)=-\log p_\theta(x)=\log(1+e^{-z(x)^\top\theta})$.
Expected reward and population cross-entropy are, respectively,
\begin{equation}
 J(\theta)=\E[p_\theta(X)],
 \qquad L(\theta)=\E[\ell(\theta,X)].
 \label{eq:classification-objectives}
\end{equation}
Each log loss is convex in $\theta$, so $L$ is convex as well.
The experiment defining $J$ draws a context $X$ from $\nu$ and then
draws the predicted label $A$ from $\pi_\theta(\cdot\mid X)$.
Because reward is one exactly when the prediction is correct,
$J(\theta)=\P(A=y(X))$: the probability of a correct prediction,
averaged over both the context and the policy's random choice of label.
Always choosing the most probable label instead gives the deterministic
classifier $h_\theta(x)=\operatorname{sign}(\theta^\top x)$, with
$\operatorname{sign}(0)=1$ as in Lecture 10. Its probability of a correct
prediction is $\P(h_\theta(X)=y(X))$, a different objective.
The population is \emph{linearly separable} if
there is a $w\in\R^d$ with $z(X)^\top w>0$ almost surely.
Then $J(cw)\to1$ as $c\to\infty$ by bounded convergence. Population
separability is stronger than observing a separable finite training sample.
Throughout the separable contextual-bandit results below,
\[
 J^\star=\sup_\theta J(\theta)=1,
 \qquad \delta(\theta)=J^\star-J(\theta)=1-J(\theta).
\]

\paragraph{Comparing the gradients.}
To isolate the effect of the objective, suppose all labels are available
and compute both population gradients exactly. At each context $x$,
the gradients of success probability and log loss are
\begin{equation}
 \begin{aligned}
 \nabla p_\theta(x)&=p_\theta(x)(1-p_\theta(x))z(x),\\
 -\nabla\ell(\theta,x)&=(1-p_\theta(x))z(x).
 \end{aligned}
 \label{eq:classification-weights}
\end{equation}
Thus $\nabla p_\theta(x)=-p_\theta(x)\nabla\ell(\theta,x)$:
at context $x$, increasing success probability gives the same direction
as decreasing log loss, scaled by $p_\theta(x)$.
When the policy assigns very small probability to the correct label at
$x$, that context has little influence on the population policy gradient.
Updates from other contexts can reduce this probability under either
objective. For log loss, convexity gives a global guarantee: on a separable
population, exact gradient flow satisfies $L(\theta(t))\to0$, as in
Lecture 11. Since $\delta(\theta)=1-J(\theta)\le L(\theta)$, we also have
$J(\theta(t))\to1$. Individual correct-label probabilities need not
increase monotonically. \Cref{prop:shared-pg} will show that policy gradient
can fail to achieve this convergence.

\subsection{When a good initialization is enough}
\label{sec:local-pg}

Under separability, maximizing $J$ is the right optimization problem:
$\sup_\theta J(\theta)=1$. Searching the entire parameter space could
therefore find policies with arbitrarily high reward. The question for
policy gradient is whether successive local improvements reach those
policies. Lecture 11 gave two sufficient ingredients: the objective must
improve along the trajectory, and its gradient must control the remaining
error. Exact gradient-ascent flow never decreases $J$, so from $\theta(0)$
it can only visit parameters satisfying $J(\theta)\ge J(\theta(0))$.
The gradient bound need only hold on that set.

\begin{proposition}[Gradient domination in a high-reward region]
\label{prop:local-pg}
Let $J:\R^d\to\R$ have a globally Lipschitz gradient and a finite supremum
$J^\star$. Put $\delta(\theta)=J^\star-J(\theta)$. Suppose that, for some
$\bar\delta,c>0$ and $r\in\{1,2\}$,
\[
 c\,\delta(\theta)^r\le\|\nabla J(\theta)\|_2^2
 \quad\text{whenever }\delta(\theta)\le\bar\delta.
\]
If $0<\delta_0=\delta(\theta(0))\le\bar\delta$, exact gradient ascent flow
$\dot\theta(t)=\nabla J(\theta(t))$ satisfies
\[
 \delta(\theta(t))\le
 \begin{cases}
 \delta_0e^{-ct},&r=1,\\
 (\delta_0^{-1}+ct)^{-1},&r=2.
 \end{cases}
\]
\end{proposition}
\begin{proof}
First check that the gradient bound remains applicable. Along the flow,
\[
 \frac{d}{dt}J(\theta(t))=\|\nabla J(\theta(t))\|_2^2\ge0,
\]
so $\delta(\theta(t))\le \delta_0\le\bar\delta$. Inside this region the error obeys
$\dot\delta\le-c\delta^r$. For $r=1$, integration gives exponential decay.
For $r=2$, $\frac{d}{dt}(1/\delta)\ge c$ while $\delta>0$, giving the reciprocal
bound. If the error reaches zero, it stays zero. These are the continuous
versions of the scalar rate arguments from Lecture 11 and
\cref{prop:drift-rate}.
\end{proof}

The proposition is a general optimization argument; the policy model must
supply its gradient bound.\endnote{Another standard local result concerns
a finite maximizer $\theta^\star$ of a $C^2$ objective with
$\nabla^2J(\theta^\star)\prec0$. In a sufficiently small ball,
$J$ is strongly concave, a gradient flow initialized in that ball remains
there, and the
local reward gap and parameter distance decay exponentially. In our
separable softmax model, reward one is approached only as parameters
diverge, so a neighborhood of a finite maximizer is unavailable.}
For the next application of the gradient bound, suppose $\nu$ is supported on
finitely many distinct points $x_1,\ldots,x_m$, with masses
$\mu_j=\nu(\{x_j\})>0$. Write $z_j=y(x_j)x_j$ and
$p_j(\theta)=\sigma(z_j^\top\theta)$, writing $p_j$ when $\theta$ is fixed, so
$J(\theta)=\sum_j\mu_jp_j(\theta)$. Along a flow, abbreviate
$p_j(t)=p_j(\theta(t))$.
The useful bound follows once every correct-label probability stays
bounded away from zero. A sufficiently high initial reward guarantees
that property.

\begin{proposition}[A sufficient initialization for policy-gradient convergence]
\label{prop:initialized-pg}
For the finite population above, let $w$ be a unit vector such that
$z_j^\top w\ge\gamma>0$ for every $j$. Define
\[
 \mu_{\min}=\min_j\mu_j,\qquad
 \delta_0=\delta(\theta(0)).
\]
If $\delta_0<\mu_{\min}$, put $\alpha=1-\delta_0/\mu_{\min}>0$.
Then exact policy-gradient flow satisfies
\[
 \delta(\theta(t))\le
 \frac{1}{\delta_0^{-1}+\gamma^2\alpha^2t}.
\]
In particular, $p_j(t)\to1$ for every $j$.
\end{proposition}
\begin{proof}
\textbf{Prevent a context from losing its correct action.}
Total error decreases along the flow, and the contribution of any one
context is at most that total:
\[
 \mu_j(1-p_j(t))\le\delta(\theta(t))\le \delta_0.
\]
Consequently $\alpha\le p_j(t)$ for every $j$ and $t$.
The initial error is too small to permit the complete loss of even the
least likely context.

\textbf{Use separation to bound the gradient.}
The separator increases every correct-label margin, so its inner product
with the gradient has no cancellation:
\[
 \gamma\alpha\delta(\theta)
 \le\gamma\sum_j\mu_jp_j(1-p_j)
 \le w^\top\nabla J(\theta)
 \le\|\nabla J(\theta)\|_2.
\]
The same calculation holds throughout the region $\delta(\theta)\le \delta_0$.
The finite logistic mixture has a globally Lipschitz gradient, since
$\|\nabla^2J(\theta)\|_{\mathrm{op}}
\le\frac14\sum_j\mu_j\|z_j\|_2^2$.
Apply \cref{prop:local-pg} with $\bar\delta=\delta_0$, $r=2$, and
$c=\gamma^2\alpha^2$.
Since $\mu_j>0$, convergence of total error to zero forces every
$p_j(t)$ to one.
\end{proof}

The condition concerns population reward, and becomes demanding when a
context has very small probability. It is sufficient, not necessary.
The next result replaces a favorable initialization by a geometric
condition and applies directly to sampled updates.

\subsection{When contexts help one another}
\label{sec:nonnegative-pg}

A successful update points along the signed feature of the sampled
context. Its effect on a second context is determined by the inner product
of their signed features. If all these inner products are nonnegative,
learning from one context preserves the progress made on every other one.

\begin{proposition}[Sampled policy gradient with nonnegative feature inner products]
\label{prop:nonnegative-pg}
In the binary contextual-bandit model, assume
$\|z(X)\|_2\le B$ almost surely, and let $w$ be a unit vector with
$w^\top z(X)\ge\gamma>0$ almost surely. Suppose
$z(x)^\top z(x')\ge0$ for $(\nu\otimes\nu)$-almost every pair.
At round $t$, draw a fresh $X_t\sim\nu$, sample
$A_t\sim\pi_{\theta_t}(\cdot\mid X_t)$, and observe
$R_t=\I\{A_t=y(X_t)\}$. For any fixed $\eta>0$, use zero-baseline REINFORCE:
\[
 \theta_{t+1}=\theta_t+\eta R_t
 \nabla\log\pi_{\theta_t}(A_t\mid X_t).
\]
For a finite deterministic initialization, define
$\alpha=\operatorname*{ess\,inf}_{x\sim\nu}p_{\theta_0}(x)>0$ and
$\delta_0=\delta(\theta_0)$. Then
\[
 \E[\delta(\theta_t)]\le
 \frac{1}{\delta_0^{-1}+\eta e^{-\eta B^2}\gamma^2\alpha^2t},
 \qquad \delta(\theta_t)\longrightarrow0\quad\text{almost surely}.
\]
\end{proposition}
\begin{proof}
\textbf{Preserve a lower bound on success probabilities.}
A failure gives no update. A success at $x$ gives the increment
$\eta(1-p_\theta(x))z(x)$. Every other correct margin increases by
\[
 h(x')=\eta(1-p_\theta(x))z(x')^\top z(x),
 \qquad 0\le h(x')\le\eta B^2.
\]
Thus $p_{\theta_t}(x')\ge\alpha$ almost surely for almost every $x'$.
Bounded features give
$\alpha\ge\sigma(-B\|\theta_0\|_2)>0$; from zero initialization,
$\alpha=1/2$.

\textbf{Relate finite-step progress to the gradient.}
We need a lower bound on the change in a probability, not just its
first-order approximation. Since
$(\log\sigma')'=1-2\sigma\ge-1$, for every scalar margin $m\in\R$
and $0\le h\le\eta B^2$,
\[
 \sigma(m+h)-\sigma(m)
 =\int_0^h\sigma'(m+s)\,ds
 \ge e^{-\eta B^2}\sigma'(m)h.
\]
Let $\cF_t$ contain all observations before round $t$.
Average this inequality over a test context and then the sampled update.
Writing $g_t=\nabla J(\theta_t)$, whose score estimate has conditional
mean $g_t$, gives
\[
 \begin{aligned}
 \E[J(\theta_{t+1})-J(\theta_t)\mid\cF_t]
 &\ge e^{-\eta B^2}
 \E[g_t^\top(\theta_{t+1}-\theta_t)\mid\cF_t]\\
 &=\eta e^{-\eta B^2}\|g_t\|_2^2.
 \end{aligned}
\]
The separator supplies the same gradient bound used above:
\[
 \gamma\alpha\delta(\theta_t)
 \le w^\top g_t\le\|g_t\|_2.
\]
For $a_t=\E[\delta(\theta_t)]$, conditional expectation and Jensen's
inequality therefore give
$a_{t+1}\le a_t-ca_t^2$, with
$c=\eta e^{-\eta B^2}\gamma^2\alpha^2$.
The reciprocal argument in \cref{prop:drift-rate} gives the stated rate.
The errors decrease on every sample path outside a null event.
Their limit has expectation zero by bounded convergence, so that limit
is zero almost surely.
\end{proof}

Orthogonal signed features satisfy the condition, but positive inner
products also allow examples to help one another. The result uses
deterministic zero--one rewards and a zero baseline; these assumptions
make every nonzero update favorable to all contexts. Separability alone
permits negative interactions and does not ensure convergence.
\Cref{app:diagonal-pg} permits some negative interactions under a weighted
diagonal-domination condition, but only for exact gradient flow initialized
with nonnegative margins.

\subsection{A separable problem that exact policy gradient fails to learn}

We want an update that increases the correct-label probability at one
context to decrease that probability at a second context. Take two equally
likely contexts with signed features
\begin{equation}
 z_1=(1,0)^\top,\qquad z_2=(-2,1)^\top,
 \qquad \theta=(u,v)^\top.
 \label{eq:two-context-features}
\end{equation}
One realization takes $x_j=z_j$ and $y(x_j)=1$ for $j=1,2$.
The margins are $u$ and $-2u+v$. Increasing $u$ improves the first margin
and decreases the second twice as fast; increasing $v$ can repair the second.
The population is separable: $w=(1,3)^\top$ gives both margins equal
to one, so scaling $w$ makes both correct-label probabilities tend to one.

\begin{proposition}[Exact policy gradient can fail on a separable population]
\label{prop:shared-pg}
For \cref{eq:two-context-features}, start at $\theta(0)=(4,0)^\top$.
Population gradient ascent flow $\dot\theta=\nabla J(\theta)$ satisfies
\[
 u(t)\to\infty,\qquad 0\le v(t)<1,
 \qquad p_1(t)\to1,\quad p_2(t)\to0.
\]
Thus $J(\theta(t))\to1/2$ and $\delta(\theta(t))\to1/2$, although
$\sup_\theta J(\theta)=1$.
From the same initialization, cross-entropy gradient flow
$\dot\theta=-\nabla L(\theta)$ satisfies $p_1(t),p_2(t)\to1$.
\end{proposition}
\begin{proof}
\textbf{Find what would prevent recovery.}
Put $a=\sigma'(u)$ and $c=\sigma'(-2u+v)$. The policy-gradient equations are
\begin{equation}
 \dot u=\frac{a-2c}{2},\qquad \dot v=\frac c2.
 \label{eq:shared-pg-flow}
\end{equation}
The factor $1/2$ comes from the context probabilities. If $v$ stays bounded
while $u$ increases, the corrective term $c$ decreases roughly as $e^{-2u}$,
whereas $a$ decreases roughly as $e^{-u}$. To prove failure, we therefore
try to bound the total increase in $v$ as $u$ grows.

\textbf{Compare correction with deterioration.}
Temporarily suppose $u\ge4$ and $0\le v\le1$. The sigmoid derivative gives
\[
 \frac14e^{-u}\le a,\qquad c\le e^{-2u+v}\le e^{1-2u}.
\]
Since $2e^{1-u}\le2e^{-3}<1/8$, the contribution $2c$ is at most
$e^{-u}/8$. Hence
\[
 \frac1{16}e^{-u}\le\dot u,
 \qquad \frac{dv}{du}=\frac{\dot v}{\dot u}\le8e^{1-u}.
\]
The ratio measures how much the corrective coordinate can change per unit
increase in $u$. Integrating from the initial value gives
\[
 0\le v(u)\le\int_4^u8e^{1-s}\,ds<8e^{-3}<1.
\]
The resulting bound verifies the region used to obtain it: $u$ increases,
$v$ is nonnegative, and $v$ cannot reach one. If $u$ had a finite upper
bound, the positive lower bound on $\dot u$ would force further growth.
Thus $u\to\infty$ and $-2u+v\to-\infty$, proving the PG claim.

\textbf{Check the supervised comparison.}
Cross-entropy $L$ is convex, and its gradient flow makes $L$ nonincreasing.
For any fixed comparator $w$, convexity gives
\[
 \frac{d}{dt}\frac{\|\theta(t)-w\|_2^2}{2}
 =-\langle\theta(t)-w,\nabla L(\theta(t))\rangle
 \le L(w)-L(\theta(t)).
\]
Integrate and use monotonicity of $L(\theta(t))$ to obtain
\[
 L(\theta(t))\le L(w)+\frac{\|\theta(0)-w\|_2^2}{2t}.
\]
For $w=b(1,3)^\top$ with $b>0$, $L(w)=\log(1+e^{-b})$.
First let $t\to\infty$, then $b\to\infty$. We obtain
$L(\theta(t))\to0$, which forces both correct-label probabilities to one.
\end{proof}

The policy-gradient failure occurs with correct rewards and exact population
gradients. The small correct-label probability at the second context allows
updates from the first context to suppress its corrective contribution further.
Fresh contexts alone do not remove this mechanism:
\cref{prop:shared-online} proves that online, zero-baseline REINFORCE with
i.i.d. contexts and any positive constant step also has positive probability
of approaching expected reward $1/2$ on this population.

\subsection{Natural gradient and representability of policy updates}

The two-context failure disappears when the feature vectors are orthogonal:
an update from one example then leaves every other margin unchanged.
For the finite population, write $m_j(\theta)=z_j^\top\theta$ for
margin $j$ and $m_j(t)=m_j(\theta(t))$ along a flow. When the current
parameter is fixed, suppress it in $m_j$ and $p_j$.

\begin{proposition}[Policy gradient learns with orthogonal features]
\label{prop:orthogonal-pg}
For a finite population with mutually orthogonal nonzero signed features
$z_1,\ldots,z_m$, exact population policy-gradient flow from any finite
initial parameter satisfies $p_j(t)\to1$ for every $j$, and hence
$\delta(\theta(t))\to0$.
\end{proposition}
\begin{proof}
Orthogonality removes all cross terms in each margin derivative:
\[
 \dot m_j=z_j^\top\nabla J(\theta)
 =\mu_j\|z_j\|_2^2\sigma'(m_j)>0.
\]
If the increasing margin $m_j$ had a finite upper bound, its derivative
would have a positive lower bound, a contradiction. Thus every margin
tends to infinity.
\end{proof}

In the two-context example, the invertible change of variables
$\boldsymbol m=\mathbf Z\theta$, where $\mathbf Z $ has rows $z_1^\top,z_2^\top$, makes the
two margins independent coordinates. The transformed features are the
coordinate vectors, so ordinary gradient flow in $\boldsymbol m$ learns both examples.
Yet ordinary gradient flow in $\theta$ fails. The family of policies is
identical; the gradient algorithm depends on its coordinates.

To see the dependence generally, write $\theta=\psi(\xi)$ for a
smooth invertible reparameterization and let $B=D\psi(\xi)$ be its
invertible Jacobian. Write $g_\theta=\nabla_\theta J(\theta)$ and
$g_\xi=\nabla_\xi(J\circ\psi)(\xi)$ for the gradients in the two
coordinate systems. The chain rule gives $g_\xi=B^\top g_\theta$.
Ordinary gradient flow in the $\xi$ coordinates therefore induces
\[
 \dot\theta=B\dot\xi=Bg_\xi=BB^\top g_\theta,
\]
which generally differs from gradient flow $\dot\theta=g_\theta$ in the
original coordinates. A rescaling of parameters can change the policy's
learning trajectory.

Natural gradient removes this dependence of the flow on coordinates.
Using Lecture 18's notation, define the \emph{policy Fisher information matrix} by
\[
 I(\theta)=\E[s_\theta(X,A)s_\theta(X,A)^\top],
 \qquad s_\theta(x,a)=\nabla\log\pi_\theta(a\mid x).
\]
When $I$ is positive definite, the \emph{natural-gradient direction}
is $h_\theta=I(\theta)^{-1}g_\theta$. Under a change of coordinates,
this direction obeys the following rule.

\begin{proposition}[Reparameterization invariance of natural-gradient flow]
\label{prop:npg-invariance}
Under the smooth invertible change of coordinates above, suppose the
Fisher matrix is positive definite. The natural directions satisfy
\[
 B h_\xi=h_\theta.
\]
Thus the flows describe the same policy trajectory from corresponding
initial conditions, while the solutions exist.
\end{proposition}
\begin{proof}
Scores transform by the same chain rule as gradients, so
$I_\xi=B^\top I_\theta B$, where $I_\theta=I(\theta)$ and $I_\xi$
is the Fisher matrix computed in the $\xi$ coordinates. Consequently,
\[
 h_\xi=(B^\top I_\theta B)^{-1}B^\top g_\theta
 =B^{-1}I_\theta^{-1}g_\theta=B^{-1}h_\theta.
\]
Multiplication by $B$ gives the stated rule and the equality of flows.
\end{proof}

The parameter direction changes coordinates; the induced policy motion
does not. This invariance is a useful property, but does not uniquely
select the Fisher matrix: any positive-definite metric transforming as
$M_\xi=B^\top M_\theta B$ gives the same covariance rule for
$M^{-1}g$. The Fisher choice comes from the policy's local KL geometry,
connecting it to mirror descent. Invariance alone gives no guarantee of
fast learning or convergence to an optimal policy.\endnote{For an increment
$h$, the context-averaged divergence satisfies
$\E[\KL(\pi_{\theta+h}(\cdot\mid X),\pi_\theta(\cdot\mid X))]
=h^\top I(\theta)h/2+o(\|h\|_2^2)$.
With $g=\nabla J(\theta)$ and $I=I(\theta)$, the maximizer of
$g^\top h-h^\top Ih/2$ is $I^{-1}g$.
This is a local quadratic approximation to KL-regularized improvement;
its steepest-direction interpretation depends on the chosen metric.}

For a fixed $\eta>0$, we use the additive NPG step
$\theta^+=\theta+\eta h_\theta$. Such finite steps are exactly
covariant under affine changes of coordinates. Under a nonlinear change,
$\psi(\xi+\eta h_\xi)=\theta+\eta h_\theta+O(\eta^2)$;
the discretization generally loses exact invariance.

Return to the $\theta$ coordinates and write $g=\nabla J(\theta)$.
For a finite population, conditional on $X=x_j$, the score is
$(\I\{A=y(x_j)\}-p_j)z_j$. Define $\mathbf Z $ by row $j$ equal to
$z_j^\top$, let $W=\operatorname{diag}(\mu_j p_j(1-p_j))$, and let
$\mathbf1\in\R^m$ be the all-ones column vector.
Then
\begin{equation}
 g=\mathbf Z ^\top W\mathbf1,\qquad I=\mathbf Z ^\top W\mathbf Z.
 \label{eq:context-fisher}
\end{equation}
All diagonal entries of $W$ are positive at finite parameters, so $I$ is positive
definite whenever the signed features span $\R^d$.

We can now give a positive convergence result. Suppose the policy class
contains a parameter direction that increases every correct-label margin
equally. Natural gradient follows precisely that direction.

\begin{proposition}[Exponential learning with a common margin increment]
\label{prop:npg-positive}
Suppose $\mathbf Z $ has full column rank and $\mathbf Z h=\mathbf1$ has a solution.
From any finite $\theta_0$, exact population NPG with fixed $\eta>0$
satisfies $\theta_t=\theta_0+t\eta h$ and
\[
 \delta(\theta_t)\le e^{-\eta t}\sum_j\mu_j e^{-z_j^\top\theta_0}.
\]
In particular, every correct-label probability tends to one.
\end{proposition}
\begin{proof}
Since $\mathbf Z h=\mathbf1$, \cref{eq:context-fisher} gives
$Ih=\mathbf Z ^\top W\mathbf Z h=\mathbf Z ^\top W\mathbf1=g$ at every parameter.
Thus each NPG step adds $\eta h$, increasing every margin by $\eta$.
The bound follows by averaging
$1-\sigma(z_j^\top\theta_0+t\eta)\le e^{-z_j^\top\theta_0-t\eta}$.
\end{proof}

The two-context example satisfies this condition because $\mathbf Z $ is invertible:
\begin{equation}
 I^{-1}g=\mathbf Z ^{-1}\mathbf1=(1,3)^\top,
 \qquad \mathbf Z \theta^+=\mathbf Z \theta+\eta\mathbf1.
 \label{eq:two-context-npg}
\end{equation}
Both correct-label log odds increase by $\eta$ on every step, and both
probabilities tend to one. In fact, the two-context class can represent
every pair of interior binary distributions: $\boldsymbol m=\mathbf Z\theta$ is an invertible
change of coordinates. NPG gives the same policy update in either set of
coordinates.

The successful NPG update has a useful policy-space interpretation. A
\emph{policy mirror-descent} step with KL regularization maximizes
\[
 \E\left[\sum_a\pi(a\mid X)r(X,a)
 -\frac1\eta\KL(\pi(\cdot\mid X),\pi_\theta(\cdot\mid X))\right]
\]
over all conditional distributions. Lecture 19's regularized optimization
calculation gives the exponential-weights update
\begin{equation}
 \pi^+(a\mid x)=
 \frac{\pi_\theta(a\mid x)e^{\eta r(x,a)}}
 {\sum_{a'}\pi_\theta(a'\mid x)e^{\eta r(x,a')}}.
 \label{eq:context-mirror}
\end{equation}
Because the correct label has reward one and the other label has reward
zero, each correct-label log odds increases by $\eta$.
For a finite population, the target is representable on the support of
$\nu$ exactly
when a parameter direction $h$ satisfies $\mathbf Z h=\mathbf1$.
Separability only requires $\mathbf Z w>0$ coordinatewise for some $w$.

These conditions differ already with three equally likely contexts:
\begin{equation}
 z_1=(1,0)^\top,\quad z_2=(0,1)^\top,\quad z_3=(-3,1)^\top.
 \label{eq:three-context-features}
\end{equation}
The vector $w=(1,4)^\top$ gives margins $(1,4,1)^\top$, so the
population is separable. Every parameter satisfies $m_3=-3m_1+m_2$;
reparameterization preserves this relation and cannot make the three
margins independent coordinates. In particular, $\mathbf Z h=\mathbf1$ requires
$h_1=h_2=1$ and $-3h_1+h_2=1$, which is impossible. Moreover, exact NPG itself
can fail on this population.

\begin{proposition}[Separation does not ensure NPG consistency]
\label{prop:shared-npg}
For \cref{eq:three-context-features}, initialize $\theta_0=(6,6)^\top$
and use exact population NPG with any fixed $0<\eta\le1/10$.
Writing $p_{j,t}=p_j(\theta_t)$, we have $p_{1,t},p_{2,t}\to1$ and
$p_{3,t}\to0$, so $J(\theta_t)\to2/3$,
and $\delta(\theta_t)\to1/3$, although $\sup_\theta J(\theta)=1$.
\end{proposition}

\Cref{app:shared-npg} proves the proposition. NPG corrects the coordinate
dependence exposed by the two-context example, but its update remains
restricted to changes the policy class can represent. Exact computation
of the natural direction does not ensure consistency under separability.

Approximating the natural direction does not require forming the full
Fisher matrix. A fresh score sample $s$ gives an unbiased matrix-vector
product $s(s^\top h)$ for $Ih$, computable with $O(d)$ arithmetic after
the score is obtained. Iterative methods can therefore approximate the
linear solve using vector operations. Their accuracy depends on sampling,
conditioning, and the computation spent on the solve; matrix approximation
and finite-step discretization are separate sources of error.

\emph{Approximate policy mirror descent} (APMD) fits successive
mirror-descent targets using sampled rewards and an inner optimization
procedure, resetting the reference to the fitted policy after each round.
Its guarantees require control of the error in approximating those targets.
For linear softmax policies, reward linearity in the policy features,
up to context-dependent constants, ensures that the targets are representable.
The three-context example violates that stronger condition. More inner
optimization cannot remove its target representation error at a fixed
reference and positive $\eta$.
\Cref{app:rebel} shows that exact population minimization of one pairwise
squared-loss objective produces precisely NPG and inherits
\cref{prop:shared-npg}. That result concerns the specified update; it does
not establish failure of every APMD implementation.

\subsection{Fitting successful responses}
\label{sec:successful-fitting}

Once a verifier has identified correct responses, we can train on those
responses with cross-entropy. When does this improve the probability of
success? We first recall the guarantee from Lecture 4, then examine two
obstacles: a policy class that cannot fit the demonstration distribution,
and training data that omit some deployment contexts.

\begin{proposition}[Reward from fitting successful responses]
\label{prop:successful-fitting}
Fix one prompt, a finite response set $\cY$, and rewards
$r:\cY\to[0,1]$. Let $q$ be a distribution supported on
$\{y:r(y)=1\}$. For a policy $\pi_\theta$, put
\[
 \begin{aligned}
 J(\theta)&=\sum_y\pi_\theta(y)r(y),
 &J^\star&=\sup_{\vartheta}J(\vartheta),\\
 L(\theta)&=-\sum_yq(y)\log\pi_\theta(y),
 &H(q)&=-\sum_yq(y)\log q(y).
 \end{aligned}
\]
Then $L(\theta)-H(q)=\KL(q,\pi_\theta)$ and
\begin{equation}
 J^\star-J(\theta)\le1-J(\theta)
 \le 1-e^{-\KL(q,\pi_\theta)}
 \le \KL(q,\pi_\theta).
 \label{eq:successful-fitting}
\end{equation}
For any parameter sequence $(\theta_t)_{t\ge0}$,
$L(\theta_t)\to H(q)$ implies $J(\theta_t)\to1$ as $t\to\infty$.
\end{proposition}
\begin{proof}
Apply \href{lecture04.pdf}{Lecture 4}, Proposition 4.2.4, to the data law
$q$, generated law $\pi_\theta$, and cost $c(y)=1-r(y)$.
The cost lies in $[0,1]$ and has expectation zero under $q$.
The log-loss identity is the one from Lecture 3.
\end{proof}

Lecture 4, Exercise 2, already asks for the bound and its sharpness.
The premise here is that excess log loss approaches zero. Minimizing loss
within a restricted policy class need not achieve that premise. Moreover,
decreasing the upper bound in \cref{eq:successful-fitting} does not imply
that the expected reward shortfall decreases.

\begin{example}[Correct demonstrations can reduce reward]
\label{ex:successful-compromise}
At one prompt, let three responses have rewards $(1,1,0)$ and use
the one-parameter policy
\[
 \pi_\theta=\frac{(e^\theta,e^{-\theta},1)}{Z(\theta)},\qquad
 Z(\theta)=e^\theta+e^{-\theta}+1.
\]
The fixed demonstration distribution is $q=(1/2,1/2,0)$: both demonstrated
responses are correct. Its population cross-entropy and the policy's
expected reward are
\[
 L(\theta)=\log Z(\theta),\qquad J(\theta)=1-\frac1{Z(\theta)}.
\]
From any $\theta(0)>0$, cross-entropy gradient flow converges to
$\theta=0$ and reduces reward to $2/3$. Exact policy-gradient flow increases
$\theta$ without bound and makes reward tend to one.
\end{example}

The policy can put nearly all its mass on either correct response.
Giving the two correct responses equal probabilities, however, forces
$\theta=0$, where the incorrect response also has probability $1/3$.
Thus the policy class can approach reward one but cannot approach $q$.
At the cross-entropy optimum,
$\KL(q,\pi_0)=\log(3/2)>0$; the first inequality in
\cref{eq:successful-fitting} is attained.
\Cref{ex:l20-successful-fitting} asks you to prove the flow claims and
check that allowing independent logits removes this conflict.

The example compares exact gradients of two objectives: cross-entropy
under fixed demonstrations and reward under the current policy.
Fresh on-policy evaluations can estimate the latter gradient, but the known
reward table already permits its exact computation. The example therefore
isolates the objectives and the policy restriction; it gives no sample
complexity separation between offline and online learning.\endnote{If
$Y\sim\pi_\theta$ and $R$ is its reward, taking the sample gradient of
$-R\log\pi_\theta(Y)$ while holding the sampled data fixed gives the
descent direction $R s_\theta(Y)$. Its expectation is $\nabla J(\theta)$
by \cref{prop:score}. Thus refreshing reward-weighted training examples
on-policy can implement PG. Fitting the fixed distribution in
\cref{ex:successful-compromise} is a different procedure.}

\begin{example}[Missing contexts can undo correct behavior]
\label{ex:missing-contexts}
Return to the binary logistic policy. The deployment population has two
equally likely contexts
$x_1=(1,0)^\top$ and $x_2=(-1,1)^\top$, both with correct label $1$.
Training contains only $x_1$. For $\theta=(u,v)^\top$, the training
objectives and deployment reward are
\[
 L_{\rm tr}(\theta)=-\log\sigma(u),\qquad
 J_{\rm tr}(\theta)=\sigma(u),\qquad
 J(\theta)=\frac{\sigma(u)+\sigma(v-u)}2.
\]
Initialize $\theta(0)=(M,2M)^\top$ for $M>0$. The deterministic classifier
is correct at both contexts, and its sampled-label reward is $\sigma(M)$,
arbitrarily close to one for large $M$.
Both cross-entropy flow $-\nabla L_{\rm tr}$ and exact PG flow
$\nabla J_{\rm tr}$ keep $v=2M$ and send $u\to\infty$.
The unseen context's correct-label probability $\sigma(2M-u)$ tends
to zero, and deployment reward tends to $1/2$.
\end{example}

The full population is separable, and the initial direction already
classifies both contexts correctly. Failure comes from continued optimization
on incomplete data: neither training objective constrains the unseen
context. Both flows eventually select the direction $(1,0)^\top$.
The example connects Lecture 10's implicit-bias analysis to the question
of identifying a rule beyond its observed inputs. It does not distinguish
PG from cross-entropy.
\Cref{ex:l20-missing-contexts} asks you to verify the dynamics.

Reusing only successful recorded actions gives a different update.
For recorded actions $Y_i$ with rewards $R_i=1$, every baseline-one score term satisfies
$(R_i-1)s_\theta(x_1,Y_i)=0$, so the parameter stays at its initial
value; multiplying those terms by finite importance weights preserves zero.
Exact PG on the training context still averages over the policy's possible
failures. At finite logits those failures have positive probability, and
fresh sampling can produce nonzero baseline-one updates.
The exercise and \cref{eq:importance} distinguish a zero update on an
all-success batch from a zero population gradient.

\section{Baselines can change future sampling}
\label{sec:baselines}

Exact-gradient comparisons do not determine the behavior of sampled
updates. A baseline preserves the conditional mean update but changes
individual increments and therefore future sampling. We return to
two-action bandits with a natural-gradient example in which one baseline
permits permanent failure and another ensures convergence.

Reward shifts expose the same distinction. Adding an action-independent
constant $c$ leaves the exact gradient unchanged because $\E[c\,s]=0$,
but changes a sampled direction from $Rs$ to $(R+c)s$. Shifting the reward
and leaving the baseline fixed can therefore change sample paths.

Take deterministic rewards $r(1)=1$ and $r(0)=1/2$, and parameterize the policy directly by a scalar
log-odds parameter $u$, with $p=\sigma(u)$ and $q=1-p$.
The expected reward is $J(u)=1/2+p/2$, with derivative $J'(u)=pq/2$.
The derivative becomes small near either deterministic policy.

The scalar \emph{Fisher information} is the second moment of the score:
\[
 I(u)=\E[(A-p)^2]=pq,\qquad A\sim\operatorname{Bernoulli}(p).
\]
The natural-gradient direction divides the reward gradient by this
quantity, giving $I(u)^{-1}J'(u)=1/2$.\endnote{For a Bernoulli policy,
$\KL(\pi_u,\pi_{u+v})=
\log(1+e^{u+v})-\log(1+e^u)-pv$.
The second derivative of $\log(1+e^u)$ is $pq$, so
$\KL(\pi_u,\pi_{u+v})=I(u)v^2/2+o(v^2)$ as $v\to0$.}
In this example the exact direction is constant; we will examine what
happens when it is estimated from one reward.

At round $t$, set $p_t=\sigma(u_t)$ and $q_t=1-p_t$, draw
$A_t\sim\operatorname{Bernoulli}(p_t)$, and observe $R_t=r(A_t)$.
For a baseline $b$ chosen before that draw, the sampled natural-gradient
direction is
\begin{equation}
 H_{b,t}=\frac{(R_t-b)(A_t-p_t)}{p_tq_t}
 =\begin{cases}(1-b)/p_t,&A_t=1,\\(b-1/2)/q_t,&A_t=0.
 \end{cases}
 \label{eq:natural-update}
\end{equation}
The update is $u_{t+1}=u_t+\eta H_{b_t,t}$. If $\cF_t$ is the history
before drawing $A_t$, then for every $\cF_t$-measurable baseline $b_t$,
$\E[H_{b_t,t}\mid\cF_t]=I(u_t)^{-1}J'(u_t)=1/2$.
Use the same fixed step $\eta>0$ and finite deterministic $u_0$ in both
baseline choices below.

\begin{proposition}[Two baselines, different limiting behavior]
\label{prop:baseline-paths}
For \cref{eq:natural-update}:
\begin{enumerate}[(a)]
\item With $b_t=0$, there is positive probability that every sampled action
is 0 and $p_t\to0$.
\item With the exact value baseline $b_t=J(u_t)$, we have $p_t\to1$
on every sample path.
\end{enumerate}
Since $J^\star=1$ and $J(u)=(1+p(u))/2$, the suboptimality gap converges
to $1/2$ in part (a) on the stated event and to zero in part (b).
\end{proposition}
\begin{proof}
\textbf{(a) Follow a possible failure trajectory.}
A sample of action 0 gives $H_{0,t}=-1/(2q_t)$. To see whether repeated selections
can prevent recovery, first compute the policy along that sequence.
Let $\bar u_t$ be the deterministic sequence obtained by always selecting 0,
and let $\bar p_t=\sigma(\bar u_t)$ and $\bar q_t=1-\bar p_t$.
Then
\[
 \bar u_{t+1}=\bar u_t-\frac{\eta}{2\bar q_t}
 \le\bar u_t-\frac\eta2,\qquad
 \bar p_t\le e^{u_0-\eta t/2}.
\]
The optimal-action probabilities are summable along this sequence.
For the event $E_N=\{A_0=\cdots=A_{N-1}=0\}$, conditional sampling gives
$\P(E_N)=\prod_{t=0}^{N-1}(1-\bar p_t)$.
To show that the infinite product is positive, choose $t_0$ so that
$\bar p_t\le1/2$ for $t\ge t_0$ and use
$-2\bar p_t\le\log(1-\bar p_t)$. Hence
\[
 -\infty<-2\sum_{t=t_0}^\infty\bar p_t
 \le\sum_{t=t_0}^\infty\log(1-\bar p_t).
\]
Every earlier factor is positive. Continuity of probability for the
decreasing events $E_N$ therefore gives
$\P(\bigcap_{N\ge1}E_N)=\prod_{t=0}^\infty(1-\bar p_t)>0$.
On this event $u_t=\bar u_t\to-\infty$, so $p_t\to0$.

\textbf{(b) Make both observations favor the better action.}
Subtracting $J(u_t)$ gives action 0 a negative advantage and action 1 a
positive advantage. Their scores have the same respective signs.
Substituting $b_t=1/2+p_t/2$ into \cref{eq:natural-update} gives
\[
 H_{J(u_t),t}=\begin{cases}q_t/(2p_t),&A_t=1,\\p_t/(2q_t),&A_t=0.
 \end{cases}
\]
Both increments are positive. If the increasing sequence $u_t$ had a finite
upper bound $U$, then $u_t\in[u_0,U]$ would keep $p_t$ and $q_t$ bounded
away from zero. Both possible values of $H_{J(u_t),t}$ would consequently have a
common positive lower bound. The update would increase $u_t$ by at least a
fixed positive amount at every step, contradicting the bound $U$.
Thus $u_t\to\infty$ and $p_t\to1$ on every path.
\end{proof}

Baseline choice therefore affects more than the variance of an estimate
at a fixed policy: individual updates alter future sampling probabilities,
and in this example change the limiting policy. The example
concerns the natural-gradient update \cref{eq:natural-update}; the
almost-sure theorem in \cref{sec:time} concerns the unpreconditioned vector
update \cref{eq:vector-pg}. The exact value baseline also uses reward means
that a learner generally must estimate. A learned baseline needs an error
analysis of its own. \Cref{ex:l20-variance-baseline} asks a separate question:
which baseline minimizes the estimator's variance at one fixed policy?

\section{Group baselines and reuse of sampled responses}
\label{sec:groups}

The preceding section used the policy's exact expected reward as a
baseline. In practice we must estimate that value. One approach is to
generate several responses to the same prompt and subtract their sample mean
reward, with the aim of reducing gradient variance. Reusing those rewards
to estimate the gradient introduces dependence: each reward contributes
to the baseline subtracted from it. We first calculate how this dependence
changes the expected update and how to correct it.
We return to the unpreconditioned score estimator of \cref{sec:feedback}.

Fix $x$ and $\theta$. Generate $n\ge2$ independent response--reward pairs
$(Y_i,R_i)$ under the experiment of \cref{sec:feedback}, and write
$s_i=s_\theta(x,Y_i)$ and $\bar R=n^{-1}\sum_{i=1}^n R_i$.
The \emph{group-centered estimator} is
\[
 g_{\rm group}=\frac1n\sum_{i=1}^n(R_i-\bar R)s_i.
\]
For each $i$, separate its own contribution from the other rewards:
\begin{align*}
 \E[\bar R s_i]
 &=\frac1n\E[R_i s_i]
   +\frac1n\sum_{j\ne i}\E[R_j]\E[s_i]\\
 &=\frac1n\nabla J_x(\theta).
\end{align*}
The cross terms vanish by independence and the zero-mean score identity.
The own-response term survives, giving the following correction.

\begin{proposition}[Group centering and leave-one-out correction]
\label{prop:group}
Under the preceding sampling experiment,
\begin{equation}
 \E[g_{\rm group}]=\left(1-\frac1n\right)\nabla J_x(\theta).
 \label{eq:group-bias}
\end{equation}
Define the \emph{leave-one-out baseline}
$b_i=(n-1)^{-1}\sum_{j\ne i}R_j$ and estimator
$g_{\rm loo}=n^{-1}\sum_i(R_i-b_i)s_i$. Then
\begin{equation}
 g_{\rm loo}=\frac{n}{n-1}g_{\rm group},\qquad
 \E[g_{\rm loo}]=\nabla J_x(\theta).
 \label{eq:loo}
\end{equation}
\end{proposition}
\begin{proof}
Subtract the preceding expression for $\E[\bar R s_i]$ from
$\E[R_i s_i]=\nabla J_x(\theta)$ and average over $i$.
For the exact correction, write
$R_i-b_i=n(R_i-\bar R)/(n-1)$ and substitute into $g_{\rm loo}$.
\end{proof}

Group centering changes the mean by a known scalar factor. Using a baseline
that excludes the current response removes that factor. Both procedures
require $n$ generated responses and $n$ evaluations per prompt.
Using rewards from several responses to the same prompt to construct
advantage estimates is one component of
\emph{group relative policy optimization} (GRPO).\endnote{%
Here the ``group'' is the sample of responses generated for the same prompt.
The advantages are ``relative'' because each reward is centered by the mean
reward of that sample; the standardized version also divides by the sample's
reward standard deviation.}\label{en:group-relative}
The identities above establish the mean update; they do not establish
a variance reduction.

\subsection{Mixed rewards and standardization}

With binary feedback of success probability $p$, all-equal groups have zero
centered advantages. Independence gives
\begin{equation}
 \P\!\left(0<\sum_{i=1}^n R_i<n\right)=1-p^n-(1-p)^n.
 \label{eq:mixed-group}
\end{equation}
The formula measures how often centering avoids an all-zero set of
advantages. It supplies no accuracy bound for the gradient estimator:
nonzero advantages can multiply cancelling scores or an unhelpful direction.
For small $p$, the inequality $1-(1-p)^n\le np$ also shows why a fixed-size
group may seldom contain a successful response.

Group-relative methods often \emph{standardize} rewards by dividing the
centered values by their empirical standard deviation. For a concrete
response-level procedure, define
\[
 v_n=\left(\frac1n\sum_{i=1}^n(R_i-\bar R)^2\right)^{1/2},
 \qquad
 \widehat A_i=\begin{cases}(R_i-\bar R)/v_n,&v_n>0,\\0,&v_n=0.
 \end{cases}
\]
The values $\widehat A_i$ are standardized group advantages.
The zero-variance convention makes the operation well-defined.
Averaging $\widehat A_i s_i$ introduces a random scaling correlated with
the group; \cref{eq:group-bias,eq:loo} no longer establish its mean.
\Cref{ex:l20-standardization} computes the resulting estimator for a binary
policy. The group calculation describes one component of GRPO; policy
ratios, clipping, and regularization further modify the update.

\subsection{Reusing responses from an older policy}

Generating and evaluating long responses can be expensive. Taking several
updates on one collected batch spreads that cost across updates, but the
sampling policy then differs from the policy being trained.
To understand the resulting correction, suppose
responses were sampled from a fixed policy $q$ with positive mass wherever
$\pi_\theta$ is positive. The evaluator law remains unchanged. Under that
sampling experiment, define
$w_\theta(Y)=\pi_\theta(Y\mid x)/q(Y\mid x)$. A direct change of measure gives
\begin{equation}
 \E[w_\theta(Y)R s_\theta(x,Y)]
 =\sum_y\pi_\theta(y\mid x)r(x,y)s_\theta(x,y)
 =\nabla J_x(\theta).
 \label{eq:importance}
\end{equation}
The parameter $\theta$ in this identity is fixed independently of the
sampled batch. If $\theta$ is fitted to that same batch, the identity does
not establish conditional unbiasedness of the empirical estimate.

\paragraph{Context-distribution mismatch.}
\Cref{eq:importance} fixes the prompt. Suppose instead that prompts belong
to a finite set, the objective averages them under a target distribution
$\nu$, and a collected batch averages them under a distribution $\mu$.
Assume $\nu(x)>0$ only if $\mu(x)>0$. Under the experiment
$X\sim\mu$, $Y\sim q(\cdot\mid X)$, followed by the usual reward draw,
assume $q(y\mid x)>0$ whenever $\pi_\theta(y\mid x)>0$. The same change of
measure gives
\begin{equation}
 \E\!\left[
 \frac{\nu(X)}{\mu(X)}
 \frac{\pi_\theta(Y\mid X)}{q(Y\mid X)}
 R s_\theta(X,Y)
 \right]
 =\nabla J_\nu(\theta),
 \quad
 J_\nu(\theta)=\sum_x\nu(x)
 \sum_y\pi_\theta(y\mid x)r(x,y).
 \label{eq:context-importance}
\end{equation}
Omitting $\nu(X)/\mu(X)$ gives the gradient of the objective averaged under
$\mu$, rather than the objective averaged under $\nu$. The response ratio
therefore does not correct a prompt-distribution mismatch.

\begin{example}[Action correction does not correct contexts]
\label{ex:offpolicy-context}
Let the contexts be $x_+$ and $x_-$, let $A\in\{0,1\}$, and set
$r(x,a)=a$. Consider the shared scalar policy
\[
 \pi_\theta(1\mid x_+)=\sigma(\theta),\qquad
 \pi_\theta(1\mid x_-)=\sigma(-\theta).
\]
For the uniform batch distribution $\mu$,
\[
 J_\mu(\theta)
 =\tfrac12\sigma(\theta)+\tfrac12\sigma(-\theta)
 =\tfrac12,
\]
so even an exactly action-corrected population policy gradient is zero.
For the target distribution $\nu$ concentrated on $x_+$,
$J_\nu(\theta)=\sigma(\theta)$ and $\nabla J_\nu(\theta)=\sigma'(\theta)>0$.
The failure comes from using the wrong context-weighted objective together
with a policy class that couples the contexts; global optimization of
$J_\mu$ cannot select the target direction either. Natural gradient computed
under $\mu$ also has zero direction. In compatible-function-approximation
terms, no scalar $v$ represents the advantage as
$v\nabla\log\pi_\theta(A\mid X)$ at both contexts: the required signs are
opposite.
\end{example}

\Needspace{4\baselineskip}
In a multi-turn problem, the policy affects the distribution of later
histories. Correcting only the action distribution conditional on a sampled
history does not correct this state-distribution mismatch; exact correction
generally requires an occupancy or trajectory ratio. Such ratios can have
high variance. \Cref{ex:l20-offpolicy-context} gives a two-instance lower
bound based on the preceding example.

For baseline one, the same change of measure gives
$\E[w_\theta(Y)(R-1)s_\theta(x,Y)]=\nabla J_x(\theta)$.
Every term is zero on an all-success batch. A batch drawn from a
full-support logging policy can have that property without contradicting unbiasedness before
conditioning on the observed rewards. Deliberately retaining only successes
changes the sampling experiment: the omitted failures prevent the same
importance-weighted formula from estimating the full gradient in general.
For autoregressive responses, $w_\theta$ is a product of token-probability
ratios. Such products can have high variance. Proximal policy optimization
(PPO) uses a clipped surrogate when reusing a collected batch.\endnote{PPO's
clipping depends on the sign of the estimated advantage. For a probability
ratio $w$, a fixed estimated advantage $\widehat A$, and $0<\varepsilon<1$,
the sample term in the clipped objective equals
\[
 \ell_\varepsilon(w,\widehat A)=
 \begin{cases}
  \widehat A\min\{w,1+\varepsilon\},&\widehat A\ge0,\\
  \widehat A\max\{w,1-\varepsilon\},&\widehat A<0.
 \end{cases}
\]
The benefit assigned to changes favored by $\widehat A$ is capped;
unfavorable changes retain the term $w\widehat A$. Hence
$\ell_\varepsilon(w,\widehat A)\le w\widehat A$: the clipped surrogate
lower-bounds the unclipped surrogate. The inequality alone supplies neither
a statistical lower confidence bound on expected reward nor a constraint
$w\in[1-\varepsilon,1+\varepsilon]$.}
Tokenwise clipping and repeated fitting to
the same batch are additional algorithmic choices, beyond the exact
response-level identity. The full practical surrogate is supplementary
reading; no convergence claim for it follows from \cref{eq:importance}.

\section{What the feedback objective guarantees}
\label{sec:accuracy}

Improving expected reward helps solve the intended task only insofar as
the reward measures success on that task. With verifiable feedback, the
test specifies the target: passing a finite test suite, for example, is
the rewarded event. If a learned reward model supplies the evaluations,
we must also relate its predictions to the intended reward. We finish
with a bound on the cost of that error and a brief offline comparison.

For a fixed prompt and strictly positive reference policy, recall the
KL-regularized objective from Lecture 19:
\[
 F_r(\pi)=\sum_y\pi(y)r(y)-\beta\KL(\pi,\pi_{\rm ref}),\qquad\beta>0.
\]
Write $\widehat r:\cY\to\R$ for the learned reward.
Let $\pi^\star$ maximize $F_r$ and $\widehat\pi$ maximize $F_{\widehat r}$
over the same policy class, assuming both maximizers exist.

\begin{proposition}[Uniform reward accuracy controls the optimized objective gap]
\label{prop:reward-error}
If $\max_y|\widehat r(y)-r(y)|\le\varepsilon$ for $\varepsilon\ge0$, then
\[
 F_r(\pi^\star)-F_r(\widehat\pi)\le2\varepsilon.
\]
\end{proposition}
\begin{proof}
Uniform reward accuracy gives $|F_r(\pi)-F_{\widehat r}(\pi)|\le\varepsilon$
for every policy. Insert the estimated objective at both policies:
\[
 F_r(\pi^\star)\le F_{\widehat r}(\pi^\star)+\varepsilon
 \le F_{\widehat r}(\widehat\pi)+\varepsilon
 \le F_r(\widehat\pi)+2\varepsilon.
\]
The middle inequality uses optimality for the estimated objective.
\end{proof}

Uniform accuracy controls the error even at responses favored by the
optimized policy. Average accuracy on logged responses gives less information
about rare responses. \emph{Offline learning} uses a fixed collected dataset;
its policy cannot request new evaluations to check such responses.
A policy $\pi$ has \emph{coverage quantity}
$C(\pi)=\sum_{y:\mu(y)>0}\pi(y)^2/\mu(y)$ relative to logging
distribution $\mu$, provided $\pi(y)=0$ whenever $\mu(y)=0$.
The quantity $C(\pi)$ grows when the policy emphasizes responses rare in the logs.
A \emph{pessimistic objective} subtracts an error allowance so that the
result is a lower bound on true reward under the assumed error control.
\Cref{app:pessimism} derives a coverage-dependent allowance and explains
how a quadratic probability-ratio penalty can implement it. The associated
constrained problem has a Lagrangian with that penalty, connecting the
statistical requirement to regularized optimization.

With online verifiable feedback, the learner can evaluate newly sampled
responses directly. The training problem then includes both finding
successful responses and updating the policy so that future successes
become more likely, as in the sampling analyses above.

\section{Changing which contexts receive updates (optional)}
\label{sec:context-weighting}

The failure example suggests a different way to use rewards: give less
weight to contexts that already succeed reliably. The next two results
study that change in the finite binary contextual-bandit model. Throughout
this section, $\mu_j>0$, $\|z_j\|_2\le B$, and a unit vector $w$
satisfies $z_j^\top w\ge\gamma>0$ for every $j$. These methods modify
the reward objective or its weighting; their guarantees do not apply to
ordinary ascent on $J$.

\subsection{Improving the worst context}

Suppose all correct-label probabilities and their gradients are available.
Define the worst-context success probability
$Q(\theta)=\min_jp_\theta(x_j)$. At a tie, improving one active context
can hurt another. To improve all active contexts, let $g(\theta)$ be the
minimum-norm vector in the convex hull of
$\{\nabla p_\theta(x_j):p_\theta(x_j)=Q(\theta)\}$.

\begin{proposition}[Worst-context reward ascent]
\label{prop:worst-context}
Along any absolutely continuous solution of
$\dot\theta(t)=g(\theta(t))$ almost everywhere, starting from a finite
parameter with $q_0=Q(\theta(0))$, we have
\[
 1-Q(\theta(t))\le
 \frac{1}{(1-q_0)^{-1}+\gamma^2q_0^2t}.
\]
Consequently, the average-reward gap satisfies
$\delta(\theta(t))\le1-Q(\theta(t))$ and obeys the same upper bound.
\end{proposition}

\Cref{ex:l20-worst-context} asks you to prove the rate and explain the
choice at ties. The direction requires evaluating all contexts and
solving a finite convex quadratic problem. The result concerns exact
gradient flow, with no sampling-cost guarantee.

\subsection{Capping updates using sampled rewards}

We can suppress updates on an already successful context without knowing
its correct label in advance. Choose $1/2\le\tau<1$, draw fresh contexts
and on-policy actions as in \cref{prop:nonnegative-pg}, and update by
\begin{equation}
 \theta_{t+1}=\theta_t+\eta R_t
 \I\{\pi_{\theta_t}(A_t\mid X_t)\le\tau\}
 \nabla\log\pi_{\theta_t}(A_t\mid X_t),\qquad\eta>0.
 \label{eq:capped-update}
\end{equation}
A rewarded action is known to be correct, so the decision to retain its
update uses only the sampled reward and the policy probability.
If another update later lowers that probability, the context becomes
eligible again. Away from threshold equalities, the mean update is the
gradient of $\sum_j\mu_j\min\{p_j,\tau\}$.

\begin{proposition}[A fixed cap and a schedule converging to optimality]
\label{prop:capped-pg}
From $\theta_0=0$, the fixed-cap algorithm makes at most
\[
 M_\tau=\frac{2\log(\tau/(1-\tau))+\eta B^2}
 {\eta(1-\tau)\gamma^2}
\]
nonzero updates. Almost surely, after finite time every $p_j>\tau$.

For convergence to reward one, suppose the finite context set is known.
Start at $\theta_0=0$ and use stages $k=1,2,\ldots$ with
$\tau_k=\sigma(k)$. Record a correct action for each context when a
reward of one is observed. After each round, check the current probability
of every recorded correct action. Advance to stage $k+1$ exactly when
all contexts have such a record and all these probabilities exceed
$\tau_k$; retain the records and parameters. Every stage finishes in
finite time almost surely, and
\[
 \min_jp_{\theta_t}(x_j)\longrightarrow1
 \quad\text{almost surely as }t\to\infty.
\]
Consequently, $\delta(\theta_t)\to0$ almost surely.
\end{proposition}

\Cref{ex:l20-capped-pg} asks you to prove both claims, including control
of the iterates between stage completions. The schedule uses no separator
or margin constant. Its certificate requires the known finite support,
stored successful actions, and evaluations of the current policy at all
contexts. A fixed cap needs no such certificate. Counting updates leaves
open the waiting time for the next eligible success; neither claim gives
a polynomial bound on the number of reward queries.

\section{Take-home messages}
\label{sec:take-home}
\begin{itemize}[leftmargin=*]
\item Policy gradients turn rewards for generated responses into parameter
updates through gradients of token log probabilities. The evaluator need
not be differentiable.
\item The cost of training depends on which responses the policy samples.
Deterministic binary feedback gives a $1/t$ expected suboptimality-gap bound and
logarithmic expected regret. A sufficiently small, horizon-independent fixed
step gives the same dependence on time for bounded stochastic rewards with
distinct means, although the proved constant can grow exponentially with the
number of actions. Larger fixed steps can give nearly linear regret despite
almost-sure convergence.
\item On a finite separable population, sufficiently high initial reward
guarantees convergence of exact policy-gradient flow. Nonnegative signed-feature
inner products guarantee convergence of sampled REINFORCE. Separability
alone permits failure despite exact rewards and gradients. Natural-gradient flow is
invariant to smooth invertible reparameterization. A common margin-increment
direction gives exponential learning; separability alone does not ensure
that such a direction exists or that NPG learns.
\item Fitting successful responses guarantees high reward when excess
log loss approaches zero. A restricted class may not fit their distribution.
Even with a separable population, optimization on data omitting contexts can
undo initially correct behavior under either cross-entropy or exact PG.
\item Baselines affect individual updates and therefore future observations.
In the natural-gradient example, preserving the mean direction does not
preserve the limiting behavior.
\item Group baselines estimate advantages from several responses to one prompt,
while batch reuse spreads generation and evaluation costs across updates. Centering,
standardization, and sampling corrections change the estimator. A response
likelihood ratio corrects reuse at a fixed prompt; a change in the prompt or
state distribution requires a separate correction. Reward accuracy and
offline coverage impose additional requirements.
\end{itemize}

\printendnotes

\section*{Bibliographical notes}
\label{sec:bibliography}

The motivation comes from the role of reinforcement learning in reasoning
model training. \citet{openai2024reasoning} report improving o1's
problem-solving performance with increased reinforcement-learning compute.
\citet{deepseek2025} use GRPO with rule-based accuracy and format rewards
for DeepSeek-R1-Zero; they report an increase in AIME 2024 average pass@1
from 15.6\% to 71.0\% during RL training. These are empirical results for
language models. The finite-policy analyses in this lecture explain
properties of the updates under their stated assumptions.

The score-function identity is the basis of REINFORCE
\citep{williams1992}. \Cref{ex:l20-unbiased-estimators} asks for an elementary
characterization under the universal-reward requirement.
\Cref{prop:gradient-gap} is the elementary
separate-logit form of the nonuniform gradient bounds used in tabular
policy-gradient analysis. \Cref{prop:drift-rate} is the standard scalar
decrease argument also used in Lecture 11's nonlinear gradient-dominance
analysis. The integral proof explains how to choose the transformation.
\Cref{thm:binary} is a self-contained classroom
specialization; no novelty is claimed for its drift calculation.
\citet{mei2024} analyze stochastic-gradient bandit convergence through
the interaction of update noise and exploration. Their expected-rate proof
factors an expectation involving the pathwise minimum optimal-arm probability
as though that random quantity were independent of the current gap.
Their Theorem~5.6 gives the inverse-initial-probability plateau lower bound
discussed after \cref{thm:fixed-step-rate}.
\citet[Appendix B.1]{baudry2025} identify the error and explain that a uniform
inverse second moment of the optimal-arm probability would repair the
argument. \Cref{thm:fixed-step-rate} supplies that moment bound and gives a
corrected expected $1/t$ rate. The permitted step is $K$ times larger than the
step prescribed by \citet{mei2024}; the corrected constant is $2M/\eta$, with
$M$ defined in \cref{eq:fixed-step-constant}. The proof of
\cref{thm:fixed-step-rate} uses a polynomial Lyapunov function that extends
the subset-product construction used by
\citet[Section 2.3]{jia2026}. \Citet{mei2025} prove the
arbitrary-constant-step almost-sure convergence result in
\cref{thm:as}. \citet[Theorem 2]{baudry2025} prove
\cref{thm:slow}; their Theorem 1 gives a complementary two-arm small-step
upper bound.

\citet{osband2026} studies allocation of a finite training budget. His
classification experiments marginalize actions exactly. The proposed
horizon loss weights an example's logit-gap gradient by
$\sigma(m+H)(1-\sigma(m))$, moving from cross-entropy toward expected
reward as the remaining horizon $H$ decreases to zero. His independent-logit
analysis and experiments motivate this modification; they do not establish
consistency with shared features. \Cref{sec:sharing} examines that separate
obstruction.

\Cref{prop:local-pg} recalls the standard gradient-domination argument
from Lecture 11. \Cref{prop:initialized-pg} verifies its assumptions for
a high-reward region of the finite separable population. Both statements
concern exact gradient flow. \Cref{prop:nonnegative-pg,prop:diagonal-pg,prop:worst-context,prop:capped-pg}
are derivations developed for these notes; their priority has not been
established. The capped-update proof uses the separator-progress and
squared-norm comparison familiar from perceptron analysis.
\citet{lin2026draft} give related noncontextual bandit counterexamples,
including failure with linearly realizable rewards, and sufficient geometric
conditions on the features. The contextual examples in
\cref{prop:shared-pg,prop:shared-online,prop:shared-npg} give self-contained
demonstrations of the failure mechanisms.

\citet{mousavi2026} study reward-weighted score updates under a behavior
policy $q$, with their main analysis omitting full importance correction.
For a fixed prompt, $K$ actions, one correct action $y^\star$, and uniform $q$,
\[
 \E[R\nabla\log\pi_\theta(Y\mid x)]
 =\frac1K\nabla\log\pi_\theta(y^\star\mid x).
\]
Thus this special case follows the supervised cross-entropy gradient in
expectation, using reward feedback to obtain its samples. Their separate
fully on-policy guarantee in Appendix A.5 permits some contexts to remain
unsuccessful; it does not establish consistency on every separable population.

\citet{li2026offlinepg} study offline policy optimization with a standalone
parameterized actor. Their local-update lower bound formalizes the
context-distribution failure in \cref{ex:offpolicy-context}: two contextual
bandit instances have the same offline population information, policy class,
rewards, and exact critics, but opposite target context distributions. Any
update using only that offline population information has average comparator
regret at least $1/4$ on one instance at every horizon, despite bounded density
ratios. The contextual-bandit data distribution in the construction is not
induced by a behavior policy under either target initial-state distribution.
Their Appendix C.3 gives an admissible MDP construction for deterministic
local updates satisfying a sign-symmetry condition. Compatible function
approximation rules out the two-context construction; without compatibility,
an NPG-style least-squares update computed under the offline distribution need
not recover the target direction.

\Cref{prop:orthogonal-pg} isolates the absence of interference.
\Cref{prop:successful-fitting} recalls Lecture 4's zero-cost forward-KL
bound, whose proof and sharpness are already in that lecture and its
Exercise 2. The elementary constructions in
\cref{ex:successful-compromise,ex:missing-contexts} distinguish
target-distribution misspecification from optimization on incomplete data;
no novelty is claimed. The second construction specializes Lecture 10's
implicit-bias discussion and also checks exact PG directly.
\Cref{prop:npg-invariance} gives the standard coordinate-transformation
argument. \Cref{prop:npg-positive} proves the positive result under the
common-increment condition and connects it to exact mirror-descent targets.
For implementation, \citet{ollivier2017} introduce TANGO: stochastic
vector updates track the Fisher solve while parameters change more slowly.
Theorem 3 proves convergence in probability to the exact natural-gradient
flow at each fixed finite time as the outer step tends to zero. The
assumptions include bounded gradient moments, a smooth Fisher matrix, and
a uniform positive lower eigenvalue bound. The result avoids explicit
Fisher estimation; it does not assert efficient approximation in every
dimension or apply uniformly along our nearly deterministic limiting
policies.

\citet{huang2026apmd} introduce the APMD framework and study concrete
instantiations. Section 2.2 assumes rewards linear in policy features;
Appendix B.4.1 characterizes target realizability up to context-dependent
constants. Section 4.3 gives the Rebel-OnPolicy objective and Appendix M
relates it to NPG. \Cref{app:rebel} derives that connection for our binary
model without reward realizability. The counterexample concerns exact
population minimization of that objective, outside the assumptions of the
paper's concrete convergence guarantees.

\citet{mei2022} study how baselines affect stochastic natural-gradient
convergence. \Cref{prop:baseline-paths} gives an elementary two-action
illustration with fully specified updates. The claim concerns that
preconditioned update, and the exact value baseline is assumed available.
\citet{shao2024} introduce GRPO in DeepSeekMath, using group-relative
advantages, tokenwise surrogates, response-length weighting, and a KL
penalty. Their experiments use learned rewards. The independent-group
identities in \cref{sec:groups} isolate operations that can also be used
with verifiable rewards. \citet{schulman2017} introduce PPO's clipped
surrogate and batch-reuse procedure.

Recent work also applies policy optimization to multi-turn language-model
agents. \citet{feng2025retool} interleave model generations with code
execution and train from task outcomes. \citet{wang2025ragen} study
trajectory-level policy optimization in interactive environments. These
systems instantiate the sequential setup of \cref{app:return}; their
algorithmic choices and empirical findings are outside the scope of these
notes.

The decreasing-step example in \cref{app:diffusion} comes from the
unpublished manuscript \emph{Lower bound on continuous-time policy gradient
with decreasing stepsizes} \citep{decreasingdraft}, consulted on
November 16, 2026. The statement concerns its Gaussian-bandit diffusion;
the full manuscript proof is outside these notes.
\citet{lattimore2026diffusion} develops a diffusion analysis with regret
upper and lower bounds. \citet[Theorem 1]{lattimore2026lyapunov} transfers
the upper-bound argument to discrete time; \cref{thm:small-step} states
that result. Its learning rate is constant during a run but depends on
the horizon and reward gaps. Section 2 of the latter paper supplies the
stopped-process and Lyapunov argument summarized in \cref{app:small-step}.

For offline preference learning, \citet{huang2024} connect pessimism to
regularization in $\chi^2$-preference optimization ($\chi$PO).
The elementary argument in \cref{app:pessimism} explains one reason for a
quadratic probability-ratio penalty. The paper uses a mixed
$\chi^2$/KL regularizer and a clipped preference loss; its guarantees
require the paper's statistical assumptions. Its $\chi^2$ convention
includes a factor $1/2$, unlike the unscaled sum in
\cref{eq:coverage-chi} below.

\section*{Glossary}
\begin{description}[style=nextline,leftmargin=1.5em,labelindent=0pt]
\item[Verifiable feedback (RLVF/RLVR)]
Computable checks supply reward evaluations of generated responses;
the evaluator need not supply a correct response to imitate.
\item[Score $s_\theta(x,y)$]
The log-probability gradient $\nabla\log\pi_\theta(y\mid x)$;
its expectation under the policy is zero.
\item[Stochastic policy gradient and REINFORCE]
A parameter update using sampled rewards multiplied by policy scores,
possibly with a baseline. This score-function update is called REINFORCE;
a preconditioner, such as inverse Fisher information, changes the direction.
\item[Value, advantage, and baseline]
The policy's value at prompt $x$ is $J_x(\theta)$; a response's advantage
is $r(x,y)-J_x(\theta)$. A value predictor estimates the expected reward.
Subtracting a baseline $b$ gives the advantage estimate $R-b$ when $b$
estimates the value. Independence from the current response, conditional on
prior information, permits baseline cancellation even if the estimate is inaccurate.
\item[Multi-turn interaction and return to go]
Model actions are interleaved with observations from a user, tool, or other
environment. An action can change the distribution of later observations.
The return to go $G_h$ contains the rewards from action $h$ onward and can
replace the full trajectory return in the score term for that action.
\item[Softmax bandit and reward laws]
Each action $a$ has a fixed reward distribution $\mu_a$ with mean $r_a$.
The policy assigns each action its own logit; sampled rewards update those
logits through the softmax score.
\item[Suboptimality gap and expected cumulative regret]
For $J^\star=\sup_\theta J(\theta)$, the suboptimality gap is
$\delta(\theta)=J^\star-J(\theta)$, with
$\delta_t=\delta(\theta_t)$ along the iterates;
$\cR_T=\sum_{t<T}\E[\delta_t]$ accumulates expected losses during training.
\item[Drift-to-rate bound]
A scalar recurrence $a_{t+1}\le a_t-ca_t^\alpha$ with $c>0$ and
$\alpha>1$ gives decay of order $t^{-1/(\alpha-1)}$.
Integrating inverse progress turns the nonlinear decrease into additive
progress; quadratic decrease gives reciprocal decay.
\item[Contextual bandit and population separation]
A context is observed before one action is chosen and rewarded. In the
classification model, $X\sim\nu$ and $z(x)=y(x)x$ is the signed feature
vector. Population separation means $z(X)^\top w>0$ almost surely for
some $w$.
\item[Capped reward updates]
Retain a rewarded score update only when the sampled action's current
probability is at most a threshold $\tau$. A finite-support stage certificate
allows thresholds to increase while preserving convergence to reward one.
\item[Fisher information and natural gradient]
The policy Fisher matrix is $I(\theta)=\E[ss^\top]$, where $s$ is the
policy score. When $I$ is invertible, the natural-gradient direction is
$I^{-1}\nabla J$. In the scalar binary example, $I(u)=p(1-p)$.
The natural direction transforms as a tangent vector under smooth invertible
changes of parameters, giving the same policy flow in either coordinate system.
\item[Policy mirror descent and target representability]
KL-regularized improvement over all policies gives
$\pi^+(a\mid x)\propto\pi(a\mid x)e^{\eta r(x,a)}$.
In the finite binary linear model, the target is representable exactly when
$\mathbf Z h=\mathbf1$ has a solution. Separability requires only $\mathbf Z w>0$.
\item[Approximate policy mirror descent (APMD)]
Successive mirror-descent targets are approximately fitted, with the fitted
policy used as the next reference. Guarantees require controlled fitting
error; separability alone does not provide that control.
\item[Successful-response fitting]
Cross-entropy fits a fixed distribution $q$ of reward-one responses.
Its excess over $H(q)$ is $\KL(q,\pi_\theta)$, which upper-bounds
the shortfall $1-J(\theta)$ from reward one. This shortfall upper-bounds
the optimization gap $J^\star-J(\theta)$. A reward-optimal policy need not
reproduce $q$.
\item[Group centering and leave-one-out baseline]
Group centering subtracts the mean of all group rewards. A leave-one-out
baseline omits the reward whose score it multiplies.
\item[Group standardization]
Division of centered rewards by a group standard deviation, with an explicit
rule for groups of zero variance.
\item[GRPO and PPO]
GRPO uses group-relative advantages in policy optimization. PPO uses a
clipped surrogate for updates based on responses from an older policy.
The estimator identities in these notes address specific components of
these procedures.
\item[Context-distribution correction]
If the target prompt distribution is $\nu$ and the batch distribution is
$\mu$, full importance correction includes the ratio $\nu(X)/\mu(X)$ in
addition to the response likelihood ratio. In a sequential problem, the
corresponding object is an occupancy or trajectory ratio.
\item[Offline pessimism and coverage]
A pessimistic objective lower-bounds true reward under an error assumption.
The coverage quantity $C(\pi)=\sum_y\pi(y)^2/\mu(y)$ controls transfer of
squared reward error from a logging distribution $\mu$ to a policy $\pi$.
\end{description}

\section{Exercises}
\label{sec:exercises}
\exerciseinstructions
\setcounter{courseexercise}{0}
\Needspace{10\baselineskip}
\begin{exercisechunk}
\item \textbf{Other unbiased gradient estimators}
\label[exercise]{ex:l20-unbiased-estimators}
\exerciseprofile{2}{1}{2}{\exproof\quad\exconstruction}{%
Identify the freedom left by requiring unbiasedness for every reward function.}
Fix a prompt $x$ and a parameter $\theta\in\R^d$.
On a finite response set $\cY$, let
$p_y=\pi_\theta(y\mid x)>0$ and
$g_y=\nabla\pi_\theta(y\mid x)\in\R^d$ for a differentiable policy.
The values $(p_y,g_y)_{y\in\cY}$ are known.
\begin{enumerate}[(a)]
\item Let $r:\cY\to\R$ be an unknown reward function. Sample
$Y\sim\pi_\theta(\cdot\mid x)$ and observe $R=r(Y)$.
Find a function $f:\cY\times\R\to\R^d$ such that
\[
 \E[f(Y,R)]=\sum_{y\in\cY}g_y r(y)
 =\nabla J_x(\theta)
 \qquad\text{for every }r:\cY\to\R,
\]
where $J_x(\theta)=\sum_y\pi_\theta(y\mid x)r(y)$.
The function $f$ may depend on the known policy probabilities and derivatives,
but must be chosen independently of $r$. Characterize all such functions,
proving both necessity and sufficiency.
\item Show that every function found in part (a) remains unbiased when
$R$ is a random reward satisfying $\E[|R|]<\infty$ and
$\E[R\mid Y=y]=r(y)$, with conditional reward laws independent of $\theta$.
\end{enumerate}
\end{exercisechunk}

\Needspace{10\baselineskip}
\begin{exercisechunk}
\item \textbf{From almost-sure convergence to sublinear expected regret.}
\label[exercise]{ex:l20-sublinear-regret}
\exerciseprofile{2}{1}{1}{\exproof}{%
Pass from convergence of the policy to vanishing expected average regret.}
Under the assumptions of \cref{thm:as}, let
\[
 \delta_t=J^\star-\sum_a\pi_t(a)r_a,\qquad
 \cR_T=\sum_{t=0}^{T-1}\E[\delta_t].
\]
Prove \cref{cor:sublinear-regret}, namely
\[
 \frac{\cR_T}{T}\longrightarrow0\qquad(T\to\infty).
\]
\end{exercisechunk}

\Needspace{10\baselineskip}
\begin{exercisechunk}
\item \textbf{REINFORCE for a binary policy.}
\label[exercise]{ex:l20-binary-update}
\exerciseprofile{1}{1}{1}{\excalculation}{%
Derive a sampled parameter update from the policy probabilities and reward.}
At a fixed input, consider the two-action softmax policy with logits
$\theta=(\theta_0,\theta_1)^\top$ and deterministic reward $r_a=a$
for $a\in\{0,1\}$. Put $u=\theta_1-\theta_0$ and $p=\sigma(u)$.
Show that $\pi_\theta(1)=p$ and
\[
 \nabla_\theta\log\pi_\theta(a)=(a-p)(-1,1)^\top.
\]
At iteration $t$, write $u_t=\theta_{t,1}-\theta_{t,0}$,
$p_t=\sigma(u_t)$, and draw
$A_t\sim\operatorname{Bernoulli}(p_t)$ conditional on the previous history.
Show that the REINFORCE update in \cref{eq:sampled-update}, with constant
step size $\eta>0$ and zero baseline, becomes
\[
 u_{t+1}=u_t+2\eta A_t(1-p_t),
\]
as stated in \cref{eq:binary-update}.
\end{exercisechunk}

\Needspace{10\baselineskip}
\begin{exercisechunk}
\item \textbf{The cost of the first successful response.}
\label[exercise]{ex:l20-first-success}
\exerciseprofile{2}{2}{2}{\exproof\quad\excalculation\quad\exinterpretation}{%
Separate a necessary sampling cost from a sufficient convergence bound.}
Consider the binary-verifier update in \cref{eq:binary-update}, with
$0<p_0<1$, $q_0=1-p_0$, and $c_0=2\eta e^{-2\eta}p_0^2$.
Let $N=\inf\{t+1:A_t=1\}$ count evaluations up to the first success.
\begin{enumerate}[(a)]
\item Show that $\P(N>t)=q_0^t$ for every integer $t\ge0$ and that
$\E[N]=1/p_0$. Deduce the lower bound
$\E[q_t]\ge q_0^{t+1}$.
\item For $0<\varepsilon<q_0$, show that $\E[q_t]\le\varepsilon$
requires
\[
 t\ge\frac{\log(q_0/\varepsilon)}{-\log q_0},
\]
whereas \cref{thm:binary} guarantees $\E[q_t]\le\varepsilon$ whenever
$t\ge(\varepsilon^{-1}-q_0^{-1})/c_0$.
For a fixed $0<\varepsilon<1$, compare the orders of these two thresholds
as $p_0\downarrow0$ with $\eta$ fixed. State whether these bounds determine
the exact order of the required number of evaluations.
\end{enumerate}
\end{exercisechunk}

\Needspace{10\baselineskip}
\begin{exercisechunk}
\item \textbf{Learning from a binary verifier with a baseline.}
\label[exercise]{ex:l20-binary-baseline}
\exerciseprofile{2}{2}{2}{\exproof\quad\excalculation}{%
Extend a convergence proof by checking which progress estimates survive a baseline.}
Use the binary logit difference and reward $r_a=a$ from \cref{sec:binary}.
Start from finite deterministic $u_0$, fix $\eta>0$ and a constant baseline
$b\in\R$, and write $p_t=\sigma(u_t)$ and $q_t=1-p_t$.
Given the history before update $t$, draw
$A_t\sim\operatorname{Bernoulli}(p_t)$ and use
\[
 u_{t+1}=u_t+2\eta(A_t-b)(A_t-p_t).
\]
\begin{enumerate}[(a)]
\item Show that $p_{t+1}\ge p_t$ for every possible action and every finite
$u_t$ if and only if $b\in[0,1]$.
\item Assume $b\in[0,1]$. Show that the conclusions of \cref{thm:binary}
still hold with the same constant $c_0=2\eta e^{-2\eta}p_0^2$:
\[
 \E[q_t]\le\frac{1}{q_0^{-1}+c_0t},\qquad
 p_t\longrightarrow1\quad\text{almost surely}.
\]
\end{enumerate}
\begin{hints}
For (b), you may use the finite-step inequality
\[
 e^{-2\eta}\sigma(u)(1-\sigma(u))\delta
 \le\sigma(u+\delta)-\sigma(u)
 \qquad(0\le\delta\le2\eta).
\]
Apply the inequality to the random logit increment, then take conditional
expectations. Once you obtain \cref{eq:binary-drift}, the remaining argument
from \cref{thm:binary} applies.
\end{hints}
\end{exercisechunk}

\Needspace{10\baselineskip}
\begin{exercisechunk}
\item \textbf{Fitting a distribution of correct responses.}
\label[exercise]{ex:l20-successful-fitting}
\exerciseprofile{2}{2}{2}{\exproof\quad\excalculation\quad\exconstruction}{%
Separate representability of a successful policy from representability of the demonstration distribution.}
In \cref{ex:successful-compromise}, the policy is
$\pi_\theta=(e^\theta,e^{-\theta},1)/Z(\theta)$,
$Z(\theta)=e^\theta+e^{-\theta}+1$, the rewards are $(1,1,0)$,
and the demonstration distribution is $q=(1/2,1/2,0)$.
\begin{enumerate}[(a)]
\item Show that $L=\log Z$, $J=1-1/Z$, and
\[
 L'(\theta)=\frac{e^\theta-e^{-\theta}}{Z(\theta)},\qquad
 J'(\theta)=\frac{e^\theta-e^{-\theta}}{Z(\theta)^2}.
\]
From any $\theta(0)>0$, prove that $\dot\theta=-L'(\theta)$
makes $\theta(t)\to0$ and $J(\theta(t))\to2/3$, whereas
$\dot\theta=J'(\theta)$ makes $\theta(t)\to\infty$ and
$J(\theta(t))\to1$.
\item Show that no parameter sequence $(\theta_t)_{t\ge0}$ satisfies
$\pi_{\theta_t}\to q$ as $t\to\infty$.
Now enlarge the class to
$\pi_{u,v}=(e^u,e^v,1)/(e^u+e^v+1)$.
Exhibit a parameter sequence along which
$L(u,v)\to H(q)=\log2$ and $J(u,v)\to1$.
\end{enumerate}
\end{exercisechunk}

\Needspace{10\baselineskip}
\begin{exercisechunk}
\item \textbf{Continued training on incomplete data.}
\label[exercise]{ex:l20-missing-contexts}
\exerciseprofile{2}{2}{2}{\exproof\quad\excalculation\quad\exinterpretation}{%
Distinguish drift on unobserved contexts from a zero sampled update on successful data.}
Use the binary logistic policy and deployment population of
\cref{ex:missing-contexts}: $x_1=(1,0)^\top$ and $x_2=(-1,1)^\top$
are equally likely and both have label $1$. Training uses only $x_1$.
Write $\theta=(u,v)^\top$ and start at $(M,2M)^\top$ for $M>0$.
\begin{enumerate}[(a)]
\item Derive the flows for $L_{\rm tr}=-\log\sigma(u)$ and
$J_{\rm tr}=\sigma(u)$, and show that in both cases
$v(t)=2M$ and $u(t)\to\infty$.
Deduce that the initial deterministic classifier is correct on both
contexts, but eventually misclassifies $x_2$.
Show that deployment reward starts at $\sigma(M)$ and tends to $1/2$,
although the deployment population is linearly separable.
\item At a fixed parameter, let $A_1,\ldots,A_n$ be independent draws
from a fixed distribution $q$ with $q(-1),q(1)>0$.
Let $R_i=\I\{A_i=1\}$ and define
\[
 \widehat g=\frac1n\sum_{i=1}^n
 \frac{\pi_\theta(A_i\mid x_1)}{q(A_i)}
 (R_i-1)\nabla\log\pi_\theta(A_i\mid x_1).
\]
Prove that $\E[\widehat g]=\nabla J_{\rm tr}(\theta)$, although
$\widehat g=0$ whenever all $R_i=1$.
Explain why the expectation identity need not hold after conditioning
on all rewards being one, and why multiplying the retained terms by
finite weights cannot restore it.
\end{enumerate}
\end{exercisechunk}

\Needspace{10\baselineskip}
\begin{exercisechunk}
\item \textbf{Which baseline minimizes variance?}
\label[exercise]{ex:l20-variance-baseline}
\exerciseprofile{2}{2}{2}{\exproof\quad\excalculation}{%
Distinguish the value baseline from the baseline minimizing score-estimator variance.}
At one fixed policy, let $S\in\R^d$ be its score and $R\in\R$ its
reward, with $\E[S]=0$, $0<\E[\|S\|_2^2]<\infty$, and
$\E[R^2\|S\|_2^2]<\infty$. For a deterministic scalar $b$, set
$g_b=(R-b)S$ and define its total variance by
$V(b)=\E[\|g_b-\E[g_b]\|_2^2]$.
\begin{enumerate}[(a)]
\item Show that the unique variance-minimizing baseline is
\[
 b^\star=\frac{\E[R\|S\|_2^2]}{\E[\|S\|_2^2]}.
\]
\item For $A\sim\operatorname{Bernoulli}(p)$ with $0<p<1$, take
$R=A$ and $S=A-p$. Show that $b^\star=1-p$ and that
$g_{b^\star}=p(1-p)$ for either action, so $V(b^\star)=0$.
Determine for which $p$ this baseline equals the value baseline $\E[R]=p$.
\end{enumerate}
\end{exercisechunk}

\Needspace{10\baselineskip}
\begin{exercisechunk}
\item \textbf{What group standardization estimates.}
\label[exercise]{ex:l20-standardization}
\exerciseprofile{2}{2}{2}{\exproof\quad\excalculation\quad\exinterpretation}{%
Identify the policy-dependent scaling introduced by a random group denominator.}
Fix $0<p<1$, and let $A_1,A_2,\ldots$ be independent Bernoulli$(p)$
variables. For a group of size $n\ge2$, set $R_i=A_i$, $S_i=A_i-p$,
$\bar A=n^{-1}\sum_{i=1}^n A_i$, and
$v_n=(n^{-1}\sum_{i=1}^n(A_i-\bar A)^2)^{1/2}$.
Define the standardized estimator by
\[
 Z_n=\begin{cases}
 \displaystyle\frac1n\sum_{i=1}^n\frac{A_i-\bar A}{v_n}S_i,&v_n>0,\\
 0,&v_n=0.
 \end{cases}
\]
\begin{enumerate}[(a)]
\item Show that $Z_n=\sqrt{\bar A(1-\bar A)}$ for every group,
including groups with all equal outcomes.
\item Show that $\E[Z_2]=p(1-p)$, but
$\lim_{n\to\infty}\E[Z_n]=\sqrt{p(1-p)}$.
You may use the strong law of large numbers and bounded convergence.
Compare the limiting expectation with the scalar reward gradient
$J'(u)=p(1-p)$, where $p=\sigma(u)$: give their ratio and explain why
no single constant, independent of $p$, corrects the limiting estimator
for every $p\in(0,1)$.
\end{enumerate}
\end{exercisechunk}

\Needspace{10\baselineskip}
\begin{exercisechunk}
\item \textbf{Off-policy context mismatch and indistinguishable targets.}
\label[exercise]{ex:l20-offpolicy-context}
\exerciseprofile{2}{2}{2}{\exproof\quad\excalculation\quad\exinterpretation}{%
Distinguish response-distribution correction from context-distribution
correction under shared policy parameters.}
Let the context set be $\{x_+,x_-\}$, the action set be $\{0,1\}$, and
$r(x,a)=a$. Let the offline context distribution $\mu$ be uniform and consider
the scalar policy class
\[
 \pi_\theta(1\mid x_+)=\sigma(\theta),\qquad
 \pi_\theta(1\mid x_-)=\sigma(-\theta).
\]
The two problem instances differ only in their target context distributions:
$\nu_+$ puts probability one on $x_+$ and $\nu_-$ puts probability one on
$x_-$. Write $J_+$ and $J_-$ for their respective expected-reward objectives.

Call a possibly randomized update rule \emph{local} for this exercise if the
conditional law of $\theta_{t+1}$, given $\theta_0,\ldots,\theta_t$, is a fixed
function of that parameter history and the joint law under $X\sim\mu$ of
\[
 \bigl(\vartheta\mapsto\pi_\vartheta(\cdot\mid X),\ r(X,\cdot)\bigr).
\]
In particular, the update rule is not given $\nu_+$ or $\nu_-$. Use the same
initialization law in the two instances, independently of the instance.
\begin{enumerate}[(a)]
\item Show that
\[
 J_+(\theta)=\sigma(\theta),\qquad
 J_-(\theta)=\sigma(-\theta),\qquad
 J_+(\theta)+J_-(\theta)=1.
\]
For $\theta_+^\star=\log 3$ and $\theta_-^\star=-\log 3$, show that
$J_+(\theta_+^\star)=J_-(\theta_-^\star)=3/4$.
\item Prove that the two runs of any local update can be coupled so that their
parameter sequences agree at every round. For $T\ge1$, define their average
comparator gaps by
\[
 \overline\delta_{T,\pm}
 =\frac1T\sum_{t=0}^{T-1}
 \E\bigl[J_\pm(\theta_\pm^\star)-J_\pm(\theta_t^\pm)\bigr],
\]
where the expectation uses the corresponding run. Show that
\[
 \overline\delta_{T,+}+\overline\delta_{T,-}=\frac12,
 \qquad
 \max\{\overline\delta_{T,+},\overline\delta_{T,-}\}\ge\frac14.
\]
\item Show that no behavior policy can induce the uniform context marginal
$\mu$ under either target distribution in this contextual bandit. Then compute
$\nu_\pm(x)/\mu(x)$ and verify
\[
 \sum_x\mu(x)\frac{\nu_\pm(x)}{\mu(x)}
 \nabla_\theta\sum_a\pi_\theta(a\mid x)r(x,a)
 =\begin{cases}
 \sigma'(\theta),&\text{for $\nu_+$},\\
 -\sigma'(\theta),&\text{for $\nu_-$}.
 \end{cases}
\]
Explain why giving the target context distribution or these context weights
violates the information restriction in the definition of a local update.
\end{enumerate}
\end{exercisechunk}

\Needspace{10\baselineskip}
\begin{exercisechunk}
\item \textbf{Average reward error after policy optimization.}\suppmark
\label[exercise]{ex:l20-offline-error}
\exerciseprofile{2}{2}{2}{\exproof\quad\excalculation\quad\exinterpretation}{%
Use coverage to convert logged reward accuracy into a policy-selection guarantee.}
Let $\cY=\{a,b\}$, with logging probabilities
$\mu(a)=1-\delta$, $\mu(b)=\delta$, where $0<\delta<1$.
The true and fitted rewards are
$r(a)=\widehat r(a)=1$, $r(b)=0$, and $\widehat r(b)=2$.
\begin{enumerate}[(a)]
\item Show that the logged squared error is $4\delta$, while maximizing
estimated expected reward over all policies selects $b$ with probability
one and incurs true-reward gap one. For that selected policy $\pi_b$,
show that $C(\pi_b)=1/\delta$ and
$\varepsilon\sqrt{C(\pi_b)}=2$, where $\varepsilon=2\sqrt\delta$.
Explain why the vanishing logged error as $\delta\downarrow0$ does not
force a vanishing policy-selection error.
\item Now consider any finite response set, logging distribution $\mu$,
and rewards satisfying, for some $\varepsilon\ge0$,
\[
 \sum_y\mu(y)(\widehat r(y)-r(y))^2\le\varepsilon^2.
\]
For $B\ge1$, let $\Pi_B=\{\pi\ll\mu:C(\pi)\le B\}$.
Let $\pi^\star$ and $\widehat\pi$ maximize true and estimated expected
reward, respectively, over $\Pi_B$.
Show that
\[
 \sum_y\pi^\star(y)r(y)-\sum_y\widehat\pi(y)r(y)
 \le2\varepsilon\sqrt B.
\]
\end{enumerate}
\begin{hints}
For part (b), apply \cref{prop:pessimism} to each policy in $\Pi_B$ and
insert the estimated reward between the two true rewards.
\end{hints}
\end{exercisechunk}

\Needspace{10\baselineskip}
\begin{exercisechunk}
\item \textbf{Improving the worst context.}\suppmark
\label[exercise]{ex:l20-worst-context}
\exerciseprofile{3}{2}{3}{\exproof\quad\excalculation\quad\exconstruction\quad\exinterpretation}{%
Choose a direction that improves every tied worst context and obtain a convergence rate.}
Consider the finite binary population from \cref{sec:context-weighting},
with signed features $z_j$, unit separator $w$, and margin $\gamma>0$.
Put $p_j(\theta)=\sigma(z_j^\top\theta)$,
$Q(\theta)=\min_jp_j(\theta)$, and
$\cA(\theta)=\{j:p_j(\theta)=Q(\theta)\}$.
Let $C(\theta)$ be the convex hull of
$\{\nabla p_j(\theta):j\in\cA(\theta)\}$, and let $g(\theta)$ be its
minimum-norm point. Consider an absolutely continuous solution of
$\dot\theta=g(\theta)$ almost everywhere; existence may be assumed.
You may use the finite-minimum chain rule: along this flow,
$\frac{d}{dt}Q(\theta(t))=\min_{j\in\cA(\theta(t))}
\nabla p_j(\theta(t))^\top g(\theta(t))$ almost everywhere.
\begin{enumerate}[(a)]
\item Prove that $u^\top g\ge\|g\|_2^2$ for every $u\in C(\theta)$,
and that equality holds for at least one active gradient.
\item Writing $q(t)=Q(\theta(t))$, show that
$\dot q(t)=\|g(\theta(t))\|_2^2$ almost everywhere and
$\gamma q(t)(1-q(t))\le\|g(\theta(t))\|_2$.
\item Prove the rate in \cref{prop:worst-context}.
Give two equally likely, linearly separable contexts and a parameter
at which an arbitrary choice of one active gradient decreases the other
active probability to first order. Explain why the minimum-norm choice
avoids that problem.
\end{enumerate}
\begin{hints}
For (a), compare $\|g+\lambda(u-g)\|_2^2$ with $\|g\|_2^2$ for
$0\le\lambda\le1$. For (c), use the two signed features from the
lecture's PG failure example at zero initialization.
\end{hints}
\end{exercisechunk}

\Needspace{10\baselineskip}
\begin{exercisechunk}
\item \textbf{Capped reward updates and increasing thresholds.}\suppmark
\label[exercise]{ex:l20-capped-pg}
\exerciseprofile{3}{3}{3}{\exproof\quad\excalculation\quad\exinterpretation}{%
Separate a bound on rewarded updates from convergence over sampling rounds.}
Use the finite binary population and capped REINFORCE update of
\cref{sec:context-weighting}. Thus fresh contexts have probabilities
$\mu_j>0$, actions are drawn from the current policy, rewards are
deterministic correctness indicators, $\|z_j\|_2\le B$, and a
unit vector $w$ satisfies $w^\top z_j\ge\gamma>0$ for every $j$.
Fix $\eta>0$ and start at $\theta_0=0$.
\begin{enumerate}[(a)]
\item For a fixed threshold $1/2\le\tau<1$, put
$m_\tau=\log(\tau/(1-\tau))$. Index the nonzero updates by $i$ and write
their increments as $c_iz_{j_i}$. Show that
$\eta(1-\tau)\le c_i\le\eta$ and that, with $S_N=\sum_{i=1}^N c_i$,
\[
 \gamma^2S_N^2\le\|\theta\|_2^2
 \le(2m_\tau+\eta B^2)S_N
\]
after $N$ such updates. Deduce the update bound in \cref{prop:capped-pg}.
\item Prove that after finite sampling time every $p_j>\tau$ almost
surely. Explain why the update bound alone is not a bound on that time.
\item Now use the increasing-threshold algorithm specified in
\cref{prop:capped-pg}, with known finite support and $\tau_k=\sigma(k)$.
Write $\theta^{(k)}$ for the parameter at the beginning of stage $k$ and
put $D=(1+\eta B^2/2)/\gamma$. Prove that every iterate within stage $k$
satisfies $\|\theta-\theta^{(k)}\|_2\le2D$ and hence
\[
 \min_jz_j^\top\theta\ge k-1-2BD.
\]
Show that every stage finishes almost surely and deduce
$\min_jp_{\theta_t}(x_j)\to1$ along all sampling rounds.
Include the possibility that the stage is waiting for a first recorded
success on a context whose current probability already exceeds its threshold.
\end{enumerate}
\begin{hints}
For (b), condition on the state immediately after any fixed number of
nonzero updates; do not condition on a purported last update as though its
time were a stopping time. For (c), compare squared distance to a suitable
translate of $\theta^{(k)}$ in the separator direction.
\end{hints}
\end{exercisechunk}


\appendix
\section{Multi-turn interaction and return to go}
\label{app:return}

In a multi-turn language-model application, a history may include user
messages, tool results, and earlier model outputs. Only actions generated by
the trained policy have policy-score terms; the laws producing the intervening
observations need not be differentiable. The trajectory score identity assigns
a terminal reward to every model action that could have affected it. When the
environment also supplies intermediate rewards, a score for the current action
need only be multiplied by rewards that can depend on that action. Past rewards
can be removed by conditioning.

Consider a finite-horizon interaction with $H$ actions, finite state and
action spaces, and rewards $R_2,\ldots,R_{H+1}$ taking values in a finite set.
As in Lecture 18, $R_{h+1}$ is observed after action $A_h$. The initial-state,
transition, and reward laws are independent of $\theta$.
Let $H_h$ be the random tuple of observations before action $A_h$,
including the current state, and sample
$A_h\sim\pi_\theta(\cdot\mid H_h)$.
Assume positive differentiable policy probabilities. Define
$s_h=\nabla\log\pi_\theta(A_h\mid H_h)$ and the
\emph{return to go} $G_h=\sum_{k=h}^H R_{k+1}$.

The policy's value before action $h$ is $\E[G_h\mid H_h]$.
The advantage of the chosen action is
\[
 \E[G_h\mid H_h,A_h]-\E[G_h\mid H_h].
\]
Subtracting the value from the observed return gives an advantage estimate:
\[
 \E[G_h-\E[G_h\mid H_h]\mid H_h,A_h]
 =\E[G_h\mid H_h,A_h]-\E[G_h\mid H_h].
\]
A value predictor can estimate the baseline from the pre-action history.
The following result permits such a baseline in the score estimator.

\begin{proposition}[Return-to-go identity]
\label{prop:return}
For the interaction just described, with objective
$J(\theta)=\E[\sum_{h=1}^H R_{h+1}]$,
\[
 \nabla J(\theta)=\E\!\left[\sum_{h=1}^H s_hG_h\right].
\]
A baseline measurable before action $h$ can be subtracted from $G_h$
without changing the expectation, provided the products are integrable.
\end{proposition}
\begin{proof}
Only policy factors in the trajectory probability depend on $\theta$.
Differentiating them gives the trajectory score $\sum_h s_h$, so
\[
 \nabla J(\theta)=\E\!\left[
 \sum_{h=1}^H s_h\sum_{k=1}^H R_{k+1}\right].
\]
Given $H_h$, every past reward $R_{k+1}$ with $k<h$ is known and
$\E[s_h\mid H_h]=0$. Hence
$\E[s_hR_{k+1}]=\E[R_{k+1}\E[s_h\mid H_h]]=0$ for $k<h$.
Removing those terms leaves $G_h$. The identical calculation applies to
any pre-action baseline.
\end{proof}

\paragraph{The Rao--Blackwell connection.}
Rao--Blackwellization is the Monte Carlo technique of replacing an estimator
by its conditional expectation, preserving its mean and reducing its
variance. To see the connection here, group the full-return estimator by
reward. The term associated with $R_{k+1}$ is
$Z_k=R_{k+1}\sum_{h=1}^H s_h$. Average over the actions sampled after that
reward is observed. Since every later score has conditional mean zero,
\[
 \overline Z_k:=\E[Z_k\mid H_{k+1}]
 =R_{k+1}\sum_{h=1}^k s_h,
 \qquad
 \sum_{k=1}^H\overline Z_k=\sum_{h=1}^H s_hG_h,
\]
where $H_{H+1}$ denotes the full history. For each reward term, the law of
total variance gives, for every $u\in\R^d$,
\[
 \operatorname{Var}(u^\top\overline Z_k)
 \leq\operatorname{Var}(u^\top Z_k).
\]
The same conditioning principle appears in
\href{lecture06.pdf}{Lecture 6}, Appendix A, where averaging over resampling
recovers the corpus mean. Here each reward term is conditioned on a
different history. The separate variance inequalities therefore do not
give a variance ordering for their sums.

The identity uses intermediate rewards only when an evaluator or environment
actually supplies them. With a single terminal reward, every $G_h$ contains
that same terminal evaluation. The identity alone supplies no finer
feedback about the correctness of intermediate reasoning. Optimizing a sum
of process rewards need not maximize the original terminal reward.

\section{A finite-horizon regret guarantee}
\label{app:small-step}

\Cref{thm:fixed-step-rate} gives a horizon-independent guarantee when all
mean rewards are distinct. The following result instead allows ties among
optimal arms. It assumes rewards in $[0,1]$, zero initialization, and a
step chosen for a prescribed finite horizon.

Consider $K\ge2$ arms. Each pull of arm $a$ gives a fresh independent reward
in $[0,1]$ with mean $r_a$. Put $J^\star=\max_a r_a$ and
\[
 \Delta_a=J^\star-r_a,\qquad
 \Delta_{\min}=\min_{a:\Delta_a>0}\Delta_a,\qquad
 \Delta_{\max}=\max_a\Delta_a.
\]
Assume at least one arm is suboptimal; optimal arms may be tied.
As in \cref{sec:bandit-setup}, $\cR_T$ denotes expected cumulative regret.

\begin{theorem}[Regret with a horizon-dependent constant step]
\label{thm:small-step}
For the bandit above, use \cref{eq:vector-pg} with $\theta_0=0$.
For every integer $T\ge K$, if
\[
 0<\eta\le\frac{\Delta_{\min}^2}
 {120\Delta_{\max}\log(TK)},
 \qquad\text{then}\qquad
 \cR_T=O\!\left(\frac{K\log K\log T}{\eta}\right),
\]
where the implicit constant is universal.
\end{theorem}

Choosing $\eta$ equal to the displayed upper bound gives
$\cR_T=O((\log T)^2)$ for a fixed bandit instance.
The step stays constant for all $T$ updates, but its choice uses $T$ and
the reward gaps. The theorem controls cumulative regret and permits tied
optimal means; it does not give a last-iterate rate for one
horizon-independent step.

\paragraph{How the proof controls rare trajectories.}
The full proof is outside these notes. With probability at least $1-1/T$,
before every update all logits exceed $-2\log(TK)$ and no optimal logit
lies more than one below a suboptimal logit.
The analysis stops the process when either condition fails and uses a
Lyapunov potential to bound expected cumulative regret before stopping.
On the exceptional event, cumulative regret is at most $T$, so its
contribution to expected regret is at most one. Quantifying that event
is essential: almost-sure convergence alone does not bound its expected cost.

\section{A decreasing-step diffusion lower bound}
\label{app:diffusion}

The lower bound below formalizes the possibility that early noise favors a
suboptimal action and later decreasing steps provide insufficient correction.
Its continuous-time model is specified here so that the result's scope can
be compared with the discrete updates in the main text.

Let $K=m+2$ and $0<\Delta\le1$. The first two arms have Gaussian reward means $1$ and
$1-\Delta$ and variance one; the remaining $m$ arms have deterministic
zero reward. Write $r=(1,1-\Delta,0,\ldots,0)^\top$ and
$\zeta=(1,1,0,\ldots,0)^\top$ for the mean and variance vectors.
Let $B_t$ be standard $K$-dimensional Brownian motion, whose coordinates
are independent. Start with $\theta_0=X_0=0\in\R^K$ and define
$\pi_{t,a}=e^{\theta_{t,a}}/\sum_b e^{\theta_{t,b}}$.
The accumulated reward process $X_t$ and logit process obey the It\^o equations
\begin{equation}
 dX_{t,a}=\pi_{t,a}r_a\,dt+\sqrt{\pi_{t,a}\zeta_a}\,dB_{t,a},
 \qquad
 d\theta_t=\eta_t(\mathrm{Id}_K-\pi_t\mathbf1^\top)\,dX_t,
 \label{eq:diffusion}
\end{equation}
where $\mathrm{Id}_K$ is the identity matrix and $\mathbf1$ is the all-one
column vector. The drift of each reward coordinate is its mean weighted by
its sampling probability, and its infinitesimal noise variance is
$\pi_{t,a}\zeta_a\,dt$. Equation \cref{eq:diffusion} defines the model;
we do not derive a diffusion limit of the discrete update here.
Here $J^\star=1$. Define the expected continuous-time regret, using the
same averaging convention as in the main text, by
\[
 \cR_T=\E\!\left[\int_0^T(J^\star-r^\top\pi_t)\,dt\right]
 =\E\!\left[\int_0^T\left(\Delta\pi_{t,2}+\sum_{a=3}^K\pi_{t,a}\right)\,dt\right].
\]

\begin{theorem}[Inverse-square-root steps in the diffusion model]
\label{thm:diffusion}
There exist absolute constants $c_m,C_m,c_\Delta,T_0,c_*>0$ such that,
for every $T\ge T_0$, integer $m\ge4$, and $\Delta>0$ satisfying
\[
 c_m\log T\le m\le C_m\frac{T}{\log T},\qquad
 \Delta\le c_\Delta\min\left\{\sqrt{\frac{\log m}{m}},\frac1{\log T}\right\},
\]
the model \cref{eq:diffusion} with $\eta_t=(1+t)^{-1/2}$ satisfies
$c_*T\Delta\le\cR_T$.
\end{theorem}

The theorem is stated without proof. For a concrete specialization,
choose $m=\lceil C\log T\rceil$ with a sufficiently large fixed $C$, and
$\Delta=c/\log T$ with a sufficiently small fixed $c>0$. For large $T$,
$1/\log T\le\sqrt{\log m/m}$ and the theorem's conditions hold.
Consequently,
\[
 \cR_T=\Omega(T/\log T),\qquad K=m+2=\Theta(\log T).
\]
The theorem uses noisy rewards and a horizon-dependent instance family.
The result shows a slow-learning mechanism for the stated diffusion, without
asserting the same lower bound for discrete-time policy gradient.

\section{Diagonal domination for exact policy gradient}
\label{app:diagonal-pg}

The nonnegative-inner-product condition in
\cref{prop:nonnegative-pg} ensures that every sampled successful update
helps every context. Exact gradient flow can permit negative interactions
because it averages the effects of all contexts before updating, provided
the initial margins are nonnegative.

Use the finite separable population from \cref{sec:local-pg}, with masses
$\mu_j$, signed features $z_j$, margins $m_j(\theta)=z_j^\top\theta$, and a
unit vector $w$ satisfying $z_j^\top w\ge\gamma>0$ for every $j$.

\begin{proposition}[Weighted diagonal domination]
\label{prop:diagonal-pg}
Suppose that, for every context $j$,
\begin{equation}
 \sum_{k:z_j^\top z_k<0}\mu_k|z_j^\top z_k|
 \le \mu_j\|z_j\|_2^2.
 \label{eq:diagonal-domination}
\end{equation}
If $m_j(\theta(0))\ge0$ for every $j$, put
$\delta_0=\delta(\theta(0))$. Exact policy-gradient flow satisfies
\[
 \delta(\theta(t))
 \le\frac{1}{\delta_0^{-1}+\gamma^2t/4}.
\]
In particular, $p_j(\theta(t))\to1$ for every $j$.
\end{proposition}
\begin{proof}
We first show that no correct-label margin becomes negative. Along the flow,
\[
 \dot m_j
 =z_j^\top\nabla J(\theta)
 =\sum_k\mu_k\sigma'(m_k)z_j^\top z_k.
\]
Suppose $m_j=0$ while every margin is nonnegative. Since
$\sigma'(m_k)\le\sigma'(0)=1/4$, discard the nonnegative cross terms and
bound the negative terms to obtain
\[
 \dot m_j
 \ge\frac14\left(
 \mu_j\|z_j\|_2^2
 -\sum_{k:z_j^\top z_k<0}\mu_k|z_j^\top z_k|
 \right)\ge0.
\]
The gradient field is locally Lipschitz, so this inward-pointing boundary
condition makes the region $m_j\ge0$ for every $j$ invariant.

Consequently, $p_j=\sigma(m_j)\ge1/2$ throughout the trajectory. The
separator now gives
\[
 \frac\gamma2\delta(\theta)
 =\frac\gamma2\sum_j\mu_j(1-p_j)
 \le\gamma\sum_j\mu_jp_j(1-p_j)
 \le w^\top\nabla J(\theta)
 \le\|\nabla J(\theta)\|_2.
\]
Thus $\|\nabla J(\theta)\|_2^2\ge\gamma^2\delta(\theta)^2/4$.
Applying \cref{prop:local-pg} proves the rate. Since every $\mu_j$ is
positive, convergence of the average gap to zero forces every $p_j$ to one.
\end{proof}

Condition \cref{eq:diagonal-domination} permits negative feature inner
products: the self-effect at each context must dominate the total weighted
harmful interaction. A single sampled update can still reduce another
margin, so the proposition does not extend the pathwise argument of
\cref{prop:nonnegative-pg}.

\section{Online sampling with shared parameters}
\label{app:shared-online}

The exact-gradient counterexample concerns a population objective. To check
whether fresh online data remove its failure, we now analyze the actual
sampled update. The key question is whether the second context can remain
unsuccessful forever while updates from the first keep reducing its success
probability.

Use the two equally likely contexts in \cref{eq:two-context-features}.
At each integer time $t\ge0$, draw $X_t\sim\nu$ independently of the
past. Given $X_t$ and the current parameter, sample
$A_t\sim\pi_{\theta_t}(\cdot\mid X_t)$ using fresh randomness, and observe
$R_t=\I\{A_t=y(X_t)\}$. Zero-baseline REINFORCE with step $\eta>0$ is
\begin{equation}
 \theta_{t+1}=\theta_t+\eta R_t(1-p_{\theta_t}(X_t))z(X_t),
 \qquad \theta_0=(4,0)^\top.
 \label{eq:shared-online}
\end{equation}

\begin{proposition}[Online REINFORCE can fail on the separable population]
\label{prop:shared-online}
For every fixed $\eta>0$, the update \cref{eq:shared-online} has positive
probability of satisfying $p_{1,t}\to1$ and $p_{2,t}\to0$.
On that event $J(\theta_t)\to1/2$ and $\delta(\theta_t)\to1/2$.
\end{proposition}
\begin{proof}
\textbf{Construct the candidate failure process.}
Take independent sequences of context draws $X_t$ and uniform random
variables $U_t,W_t$ on $[0,1]$, with all three sequences mutually
independent. Use $U_t$ to sample labels on context 1 and $W_t$ on context 2.
Consider an auxiliary process that suppresses every update on context 2.
Its parameter has the form $(\bar u_t,0)^\top$. On context 1 it increases
$\bar u_t$ by $\eta(1-\sigma(\bar u_t))$ when
$U_t\le\sigma(\bar u_t)$, and otherwise makes no update.

\textbf{Quantify how quickly the second probability decreases.}
Index successful context-1 updates by $k$, and let $b_k$ be $e^{\bar u}$
after $k$ such updates. Then
\[
 b_{k+1}=b_k\exp\left(\frac{\eta}{1+b_k}\right),
 \qquad b_0=e^4.
\]
The elementary inequality $s\le e^s-1$ gives
\[
 \eta\sigma(4)\le\frac{\eta b_k}{1+b_k}
 \le b_{k+1}-b_k,
 \qquad e^4+k\eta\sigma(4)\le b_k.
\]
Let $N_t$ count successful auxiliary updates before time $t$.
Since $\bar u_t\ge4$, every event
$\{X_t=x_1,\ U_t\le\sigma(4)\}$ produces such an update. These events
are independent with probability $\sigma(4)/2$, so the strong law gives
$ct\le N_t$ eventually almost surely for some fixed $c>0$.
The auxiliary probability of success on context 2 therefore satisfies
\[
 \bar p_{2,t}=\sigma(-2\bar u_t)
 \le\frac1{(e^4+\eta\sigma(4)N_t)^2},
 \qquad \sum_{t=0}^\infty\bar p_{2,t}<\infty
 \quad\text{almost surely}.
\]

\textbf{Show that the actual learner can follow this process forever.}
Condition on the entire context sequence and all $U_t$; call this
information $\cG$. The auxiliary trajectory is then fixed, while the
$W_t$ remain independent uniform variables. Hence the event
$E=\{W_t>\bar p_{2,t}\text{ for every }t\text{ with }X_t=x_2\}$ satisfies
\[
 \P(E\mid\cG)=\prod_{t:X_t=x_2}(1-\bar p_{2,t})>0
 \quad\text{almost surely}.
\]
Each factor is positive and the subtracted probabilities are summable;
the logarithm argument in the proof of \cref{prop:baseline-paths} proves
positivity of the infinite product. Averaging gives $\P(E)>0$.
On $E$, the actual and auxiliary processes coincide by induction, because
context 2 always gives reward zero. Thus $v_t=0$, $u_t\to\infty$, and
the two correct-label probabilities tend to one and zero, respectively.
\end{proof}

The failure event occurs despite infinitely many visits to each context.
The zero-baseline update makes no use of an incorrect label's zero reward,
while successful updates on the first context reduce the probability of
ever obtaining a positive reward on the second.

\section{A separable population on which exact NPG fails}
\label{app:shared-npg}

The main text distinguished separation from representability of a
mirror-descent update. We now prove that the three-context example also
defeats exact natural-gradient steps themselves. The proof tracks whether
the second coordinate can grow enough faster than the first to repair the
third margin.

\begin{proposition}[Restatement of \cref{prop:shared-npg}]
For the three equally likely contexts in \cref{eq:three-context-features},
initialize $\theta_0=(6,6)^\top$ and use exact population NPG with any
fixed $0<\eta\le1/10$. Then $p_{1,t},p_{2,t}\to1$ and $p_{3,t}\to0$,
so $J(\theta_t)\to2/3$ and $\delta(\theta_t)\to1/3$, although
$\sup_\theta J(\theta)=1$.
\end{proposition}
\begin{proof}
\textbf{Compute the natural direction.}
At $\theta=(u,v)^\top$, put
$a=\sigma'(u)$, $b=\sigma'(v)$, and $c=\sigma'(-3u+v)$.
\Cref{eq:context-fisher} gives
\[
 I=\frac13\begin{pmatrix}a+9c&-3c\\-3c&b+c\end{pmatrix},
 \qquad g=\frac13\begin{pmatrix}a-3c\\b+c\end{pmatrix}.
\]
The first two feature vectors span $\R^2$, so $I$ is invertible.
Define $\alpha=9c/a$ and $\beta=c/b$. Solving $Ih=g$ yields
\[
 h_u=\frac{1+\beta}{1+\alpha+\beta},
 \qquad h_v=1+\frac{3\beta}{1+\alpha+\beta}.
\]
When $c$ is small relative to $a$ and $b$, this direction is close to
$(1,1)^\top$, which decreases the third margin.

\textbf{Find a quantity that would have to grow for recovery.}
Let $d=v-u$. The third margin is $-2u+d$.
For the third margin to become positive, $d$ must exceed $2u$.
For one NPG step, write $\Delta u=\eta h_u$ and
$\Delta d=\eta(h_v-h_u)$. Then
\[
 0<\Delta u\le\eta,
 \qquad
 0\le\frac{\Delta d}{\Delta u}
 =\frac{\alpha+3\beta}{1+\beta}.
\]
Suppose for now that $u\ge6$ and $0\le d\le1$. The same
sigmoid-derivative bounds as in \cref{prop:shared-pg} give
\[
 \alpha\le36e^{1-u},\qquad \beta\le4e^{2-u},
 \qquad \Delta d\le C e^{-u}\Delta u,
 \quad C=36e+12e^2.
\]

\textbf{Bound the total increase of that quantity.}
Because $\Delta u\le\eta$, for $s\in[u,u+\Delta u]$ we have
$e^{-u}\le e^\eta e^{-s}$. Thus
\[
 \Delta d\le C e^\eta\int_u^{u+\Delta u}e^{-s}\,ds.
\]
Starting with $u_0=6$ and $d_0=0$, summing over steps gives
\[
 0\le d_t\le C e^\eta\int_6^{u_t}e^{-s}\,ds
 <C e^{\eta-6}\le C e^{-5.9}<0.511<1.
\]
Induction now verifies the region used in the bound. Moreover, throughout
that region,
\[
 \frac{\eta}{1+(36e+4e^2)e^{-6}}\le\Delta u.
\]
Therefore $u_t\to\infty$, $v_t=u_t+d_t\to\infty$, and
$-3u_t+v_t=-2u_t+d_t\to-\infty$.
The limiting correct-label probabilities are consequently $(1,1,0)$.
\end{proof}

\section{Pairwise squared-loss fitting and NPG}
\label{app:rebel}

A method can solve its inner fitting problem exactly while failing to
approach an optimal policy. We verify this distinction for a pairwise
squared-loss update, whose fitting residual measures errors in the desired
changes of log odds.

Fix a reference parameter $\theta$ in the finite binary classification
model. Draw $X\sim\nu$, then draw labels $A,A'$ independently from
$\pi_\theta(\cdot\mid X)$. Observe their rewards
$r(X,A),r(X,A')$. For a candidate increment $\eta h$, define the residual
\[
 f_h(x,a)=\frac1\eta\log\frac{\pi_{\theta+\eta h}(a\mid x)}
 {\pi_\theta(a\mid x)}-r(x,a),
 \qquad L_\theta(h)=\E[(f_h(X,A)-f_h(X,A'))^2].
\]
An idealized population update minimizes $L_\theta$ exactly and then resets
the reference to $\theta+\eta h$. Sample-based implementations replace the
expectation by a batch average and take finitely many optimization steps.

\begin{proposition}[Exact pairwise fitting gives the NPG update]
\label{prop:pairwise-npg}
If the signed features span $\R^d$, then
\[
 L_\theta(h)=2\sum_j\mu_j p_j(1-p_j)(z_j^\top h-1)^2,
\]
and the unique minimizer is $h^\star=I(\theta)^{-1}\nabla J(\theta)$.
Thus exact population minimization on the three-context population,
initialized at $(6,6)^\top$ with $0<\eta\le1/10$, approaches expected reward
$2/3$ and suboptimality gap $1/3$, even though the population is separable.
\end{proposition}
\begin{proof}
Conditional on $X=x_j$, write $B=\I\{A=y(x_j)\}$ and
$B'=\I\{A'=y(x_j)\}$. The difference of residuals is
\[
 f_h(x_j,A)-f_h(x_j,A')=(B-B')(z_j^\top h-1).
\]
The pair has unequal labels with probability $2p_j(1-p_j)$, proving
the quadratic formula. Differentiating gives
$\nabla_h L_\theta(h)=4(Ih-g)$ with $I,g$ from
\cref{eq:context-fisher}. Positive definiteness of $I$ makes $I^{-1}g$
the unique minimizer. Applying \cref{prop:shared-npg} proves the final claim.
\end{proof}

The minimum fitting loss is zero exactly when $\mathbf Z h=\mathbf1$ has a
solution. In the three-context example the minimum is positive at every
finite reference parameter, although the convex quadratic can be minimized
exactly. The result concerns this population update. It neither establishes
failure of every finite inner-loop implementation nor contradicts an APMD
guarantee conditional on sufficiently accurate target fitting.

\section{Offline reward error and pessimism}
\label{app:pessimism}

Logged data typically control error under their own sampling distribution.
A policy chosen by optimizing the fitted reward changes that distribution.
We need a bound that quantifies the cost of the change before treating the
fitted reward as a reliable objective.

Let $\varepsilon\ge0$ and let $\mu$ be a logging distribution on the finite
response set. Suppose
$e(y)=\widehat r(y)-r(y)$ satisfies
$\sum_y\mu(y)e(y)^2\le\varepsilon^2$.
Restrict attention to policies $\pi\ll\mu$, meaning that
$\mu(y)=0$ implies $\pi(y)=0$. Define the \emph{coverage quantity}
\[
 C(\pi)=\sum_{y:\mu(y)>0}\frac{\pi(y)^2}{\mu(y)}.
\]
Large $C(\pi)$ allows the policy to put substantial weight on responses
that are rare in the logs.

\begin{proposition}[A regularized lower bound on reward]
\label{prop:pessimism}
Under this squared-error assumption, for every $\pi\ll\mu$ and $\lambda>0$,
\begin{equation}
 \left|\sum_y\pi(y)e(y)\right|
 \le\varepsilon\sqrt{C(\pi)}
 \le\lambda C(\pi)+\frac{\varepsilon^2}{4\lambda}.
 \label{eq:offline-error}
\end{equation}
Consequently,
\begin{equation}
 \sum_y\pi(y)\widehat r(y)-\lambda C(\pi)
       -\frac{\varepsilon^2}{4\lambda}
 \le\sum_y\pi(y)r(y).
 \label{eq:pessimistic-objective}
\end{equation}
\end{proposition}
\begin{proof}
To use the known error average, express the desired error under the logging
distribution with a probability-ratio weight. Cauchy--Schwarz gives
\[
 \left|\sum_y\mu(y)\frac{\pi(y)}{\mu(y)}e(y)\right|
 \le\left(\sum_y\frac{\pi(y)^2}{\mu(y)}\right)^{1/2}
       \left(\sum_y\mu(y)e(y)^2\right)^{1/2}
 \le\varepsilon\sqrt{C(\pi)},
\]
where the sums are over $\mu(y)>0$.
For the second bound, expand
$0\le(\sqrt{\lambda C(\pi)}-\varepsilon/(2\sqrt\lambda))^2$.
Finally, subtract the error allowance from estimated reward to obtain
\cref{eq:pessimistic-objective}.
\end{proof}

For a fixed $\lambda$, the term $\varepsilon^2/(4\lambda)$ is constant
across policies. Maximizing the lower bound therefore amounts to maximizing
estimated reward minus $\lambda C(\pi)$.
The penalty is also the policy-dependent part of the Lagrangian for
\[
 \max_{\pi\ll\mu}\sum_y\pi(y)\widehat r(y)
 \quad\text{subject to}\quad C(\pi)\le B,
\]
whose Lagrangian is
$\sum_y\pi(y)\widehat r(y)-\lambda(C(\pi)-B)$, with $\lambda\ge0$.
The identity
\begin{equation}
 C(\pi)-1=\sum_{y:\mu(y)>0}\mu(y)
 \left(\frac{\pi(y)}{\mu(y)}-1\right)^2
 \label{eq:coverage-chi}
\end{equation}
identifies the penalty, up to an additive constant, with a chi-squared
divergence. This is a concrete way a regularized objective can implement
pessimism under a stated error bound. \Cref{ex:l20-offline-error} shows why
average accuracy alone allows poor policy selection and how a coverage
constraint restores a guarantee.

\begin{thebibliography}{99}
\bibitem[OpenAI(2024)]{openai2024reasoning}
OpenAI.
\newblock Learning to reason with LLMs.
\newblock September 12, 2024.
\newblock \url{https://openai.com/index/learning-to-reason-with-llms/}.

\bibitem[DeepSeek-AI(2025)]{deepseek2025}
DeepSeek-AI.
\newblock DeepSeek-R1: Incentivizing reasoning capability in LLMs via
reinforcement learning.
\newblock arXiv:2501.12948v1, 2025.
\newblock \url{https://arxiv.org/abs/2501.12948v1}.

\bibitem[Williams(1992)]{williams1992}
Ronald J. Williams.
\newblock Simple statistical gradient-following algorithms for connectionist
reinforcement learning.
\newblock \emph{Machine Learning}, 8:229--256, 1992.
\newblock \url{https://doi.org/10.1007/BF00992696}.

\bibitem[Mei et~al.(2023)]{mei2024}
Jincheng Mei, Zixin Zhong, Bo Dai, Alekh Agarwal, Csaba Szepesv\'ari,
and Dale Schuurmans.
\newblock Stochastic gradient succeeds for bandits.
\newblock \emph{Proceedings of ICML}, 2023.
\newblock \href{https://proceedings.mlr.press/v202/mei23a/mei23a.pdf}{Published paper};
arXiv:2402.17235.

\bibitem[Mei et~al.(2024)]{mei2025}
Jincheng Mei, Bo Dai, Alekh Agarwal, Sharan Vaswani, Anant Raj,
Csaba Szepesv\'ari, and Dale Schuurmans.
\newblock Small steps no more: Global convergence of stochastic gradient
bandits for arbitrary learning rates.
\newblock NeurIPS, 2024; updated arXiv version, 2025.
\newblock \url{https://arxiv.org/abs/2502.07141}.

\bibitem[Baudry et~al.(2025)]{baudry2025}
Dorian Baudry, Emmeran Johnson, Simon Vary, Ciara Pike-Burke,
and Patrick Rebeschini.
\newblock Does stochastic gradient really succeed for bandits?
\newblock NeurIPS, 2025.
\newblock \href{https://proceedings.neurips.cc/paper_files/paper/2025/file/a4e683f0ce6b91e7fbdae9d32642d88f-Paper-Conference.pdf}{Conference paper}.

\bibitem[Mei et~al.(2022)]{mei2022}
Jincheng Mei, Wesley Chung, Valentin Thomas, Bo Dai,
Csaba Szepesv\'ari, and Dale Schuurmans.
\newblock The role of baselines in policy gradient optimization.
\newblock NeurIPS, 2022.
\newblock \href{https://papers.neurips.cc/paper_files/paper/2022/file/718d02a76d69686a36eccc8cde3e6a41-Paper-Conference.pdf}{Conference paper}.

\bibitem[Osband(2026)]{osband2026}
Ian Osband.
\newblock Planning to learn.
\newblock arXiv:2610.03667v1, 2026.
\newblock \url{https://arxiv.org/abs/2610.03667v1}.

\bibitem[Ollivier(2017)]{ollivier2017}
Yann Ollivier.
\newblock True asymptotic natural gradient optimization.
\newblock arXiv:1712.08449v1, 2017.
\newblock \url{https://arxiv.org/abs/1712.08449v1}.

\bibitem[Lin et~al.(2026)]{lin2026draft}
Max Qiushi Lin, Jincheng Mei, Matin Aghaei, Michael Lu, Bo Dai,
Alekh Agarwal, Dale Schuurmans, Csaba Szepesv\'ari, and Sharan Vaswani.
\newblock Rethinking the global convergence of softmax policy gradient with
linear function approximation: The case of multi-armed bandits.
\newblock Unpublished manuscript under submission; second-submission draft
consulted November 16, 2026.

\bibitem[Mousavi-Hosseini and Erdogdu(2026)]{mousavi2026}
Alireza Mousavi-Hosseini and Murat A. Erdogdu.
\newblock Post-training with policy gradients: Optimality and the base model barrier.
\newblock arXiv:2603.06957v1, 2026.
\newblock \url{https://arxiv.org/abs/2603.06957v1}.

\bibitem[Li et~al.(2026)]{li2026offlinepg}
Xiang Li, Yuheng Zhang, and Nan Jiang.
\newblock Beyond state-wise mirror descent: Offline policy optimization with
parametric policies.
\newblock arXiv:2602.23811, 2026.
\newblock \url{https://arxiv.org/abs/2602.23811}.

\bibitem[Huang et~al.(2026)]{huang2026apmd}
Audrey Huang, Nan Jiang, Sham M. Kakade, Sadhika Malladi,
Akshay Krishnamurthy, and Dylan J. Foster.
\newblock A theory of policy optimization in post-training: Coverage collapse
and reference resets.
\newblock arXiv:2611.14467v1, 2026.
\newblock \url{https://arxiv.org/abs/2611.14467v1}.

\bibitem[Shao et~al.(2024)]{shao2024}
Zhihong Shao, Peiyi Wang, Qihao Zhu, Runxin Xu, Junxiao Song, Xiao Bi,
Haowei Zhang, Mingchuan Zhang, Y. K. Li, Y. Wu, and Daya Guo.
\newblock DeepSeekMath: Pushing the limits of mathematical reasoning in
open language models.
\newblock arXiv:2402.03300, 2024.
\newblock \url{https://arxiv.org/abs/2402.03300}.

\bibitem[Schulman et~al.(2017)]{schulman2017}
John Schulman, Filip Wolski, Prafulla Dhariwal, Alec Radford, and Oleg Klimov.
\newblock Proximal policy optimization algorithms.
\newblock arXiv:1707.06347, 2017.
\newblock \url{https://arxiv.org/abs/1707.06347}.

\bibitem[Feng et~al.(2025)]{feng2025retool}
Jiazhan Feng, Shijue Huang, Xingwei Qu, Ge Zhang, Yujia Qin,
Baoquan Zhong, Chengquan Jiang, Jinxin Chi, and Wanjun Zhong.
\newblock ReTool: Reinforcement learning for strategic tool use in LLMs.
\newblock arXiv:2504.11536, 2025.
\newblock \url{https://arxiv.org/abs/2504.11536}.

\bibitem[Wang et~al.(2025)]{wang2025ragen}
Zihan Wang, Kangrui Wang, Qineng Wang, Pingyue Zhang, Linjie Li,
Zhengyuan Yang, Kefan Yu, Minh Nhat Nguyen, Licheng Liu, Eli Gottlieb,
Monica Lam, Yiping Lu, Kyunghyun Cho, Jiajun Wu, Li Fei-Fei,
Lijuan Wang, Yejin Choi, and Manling Li.
\newblock RAGEN: Understanding self-evolution in LLM agents via multi-turn
reinforcement learning.
\newblock arXiv:2504.20073, 2025.
\newblock \url{https://arxiv.org/abs/2504.20073}.

\bibitem[Unpublished manuscript(2026)]{decreasingdraft}
\emph{Lower bound on continuous-time policy gradient with decreasing stepsizes}.
\newblock Unpublished manuscript, draft consulted November 16, 2026.

\bibitem[Lattimore(2026a)]{lattimore2026diffusion}
Tor Lattimore.
\newblock A diffusion analysis of policy gradient for stochastic bandits.
\newblock arXiv:2603.10219, 2026.
\newblock \url{https://arxiv.org/abs/2603.10219}.

\bibitem[Lattimore(2026b)]{lattimore2026lyapunov}
Tor Lattimore.
\newblock A Lyapunov analysis of softmax policy gradient for stochastic bandits.
\newblock arXiv:2603.26547, 2026.
\newblock \url{https://arxiv.org/abs/2603.26547}.

\bibitem[Jia and Ouyang(2026)]{jia2026}
Yuling Jia and Dong Ouyang.
\newblock Convergence and regret of the policy gradient for multi-armed
bandits in diffusion environment.
\newblock arXiv:2607.29593v1, 2026.
\newblock \url{https://arxiv.org/abs/2607.29593v1}.

\bibitem[Huang et~al.(2024)]{huang2024}
Audrey Huang, Wenhao Zhan, Tengyang Xie, Jason D. Lee, Wen Sun,
Akshay Krishnamurthy, and Dylan J. Foster.
\newblock Correcting the mythos of KL-regularization: Direct alignment
without overoptimization via chi-squared preference optimization.
\newblock arXiv:2407.13399, 2024; revised 2025.
\newblock \url{https://arxiv.org/abs/2407.13399}.
\end{thebibliography}
\end{document}
